\documentclass[11pt]{article}

\usepackage[margin=1.1in]{geometry}
\usepackage{amsmath,amssymb,amsthm,mathtools}
\usepackage{microtype}
\usepackage{booktabs}
\usepackage{array}
\usepackage{enumitem}
\usepackage[hidelinks]{hyperref}

\newtheorem{theorem}{Theorem}[section]
\newtheorem{proposition}[theorem]{Proposition}
\newtheorem{lemma}[theorem]{Lemma}
\newtheorem{corollary}[theorem]{Corollary}
\theoremstyle{definition}
\newtheorem{definition}[theorem]{Definition}
\theoremstyle{remark}
\newtheorem{remark}[theorem]{Remark}
\newtheorem{question}[theorem]{Question}

\numberwithin{equation}{section}

\DeclareMathOperator{\rk}{rk}
\DeclareMathOperator{\srk}{srk}
\DeclareMathOperator{\rank}{rank}
\DeclareMathOperator{\tr}{tr}
\DeclareMathOperator{\ld}{ld}
\DeclareMathOperator{\pr}{pr}
\DeclareMathOperator{\diag}{diag}
\DeclareMathOperator{\im}{im}
\DeclareMathOperator{\ran}{ran}
\DeclareMathOperator{\cnt}{cnt}
\DeclareMathOperator{\GL}{GL}
\DeclareMathOperator{\spn}{sp}
\DeclareMathOperator{\RePart}{Re}

\newcommand{\RR}{\mathrm{RR}}
\newcommand{\RUF}[3]{\prod_{#1}^{\mathrm{rk}} M_{#2}(#3)}
\newcommand{\Q}{\mathbb{Q}}
\newcommand{\Z}{\mathbb{Z}}
\newcommand{\R}{\mathbb{R}}
\newcommand{\C}{\mathbb{C}}
\newcommand{\N}{\mathbb{N}}
\newcommand{\F}{\mathbb{F}}
\newcommand{\U}{\mathcal{U}}
\newcommand{\HH}{\mathcal{H}}
\newcommand{\PP}{\mathcal{P}}
\newcommand{\LL}{\mathcal{L}}
\newcommand{\TT}{\mathcal{T}}
\newcommand{\NN}{\mathcal{N}}
\newcommand{\OO}{\mathcal{O}}
\newcommand{\1}{\mathbf{1}}
\newcommand{\code}[1]{{\normalfont\small\texttt{#1}}}

\title{Regular Rank \texorpdfstring{$\Q$}{Q}-Algebras Without Homomorphisms into Characteristic-Zero Matrix Ultraproducts}
\author{Jack Ryan}
\date{November 2026}

\begin{document}
\maketitle

\begin{abstract}
We construct a countable von Neumann regular $\Q$-algebra $S$ with a faithful rank function $\rk\colon S \to [0,1]$ admitting no unital ring homomorphism into any rank ultraproduct of matrix rings $M_{d_i}(F_i)$ ($d_i \ge 1$) over characteristic-zero fields $F_i$, and in particular none into Elek's complex matrix rank ultraproducts $M^{\mathrm{alg}}_\nu$. We also construct a regular rank $\C$-algebra with the same obstruction, and exhibit a finite system $B_7^\sigma$ of integer polynomial relations indexed by $\mathbb{P}^2(\F_7)$ satisfied in $S$ and in $M_{3^{3k+1}}(\F_7)$ ($k \ge 0$), yet admitting no $\varepsilon$-approximate solution in $M_N(F)$ for a uniform $\varepsilon > 0$, all characteristic-zero fields $F$, and all $N \ge 1$. This answers a reading of Elek's Section~1.2 question (arXiv:1006.5378) negatively and yields stably finite non-linearly-sofic algebras over $\Q$ and $\C$ with a faithful Sylvester matrix rank function, contrasting with Greenfeld's examples (arXiv:2210.11650). The proof pairs a characteristic-zero incidence rank bound for projective-plane idempotents with an ultrametric regular rank ring over $\Q_7$ whose $7$-adic flag tensor powers break the bound at $q=7$ since $63 < 64$. The unital non-embedding results (Theorem~\ref{thm:main-ultra}, and Theorem~\ref{thm:finite-form} for a fixed bijection $\sigma$; see Section~\ref{sec:formal}) are formally verified in Lean~4. By paper-only arguments not formalized in Lean~4, we establish the stable-finiteness contrast (Proposition~\ref{prop:stably-finite}) and extend the incidence bound to the affiliated-operator ring $\U(M)$ of any finite von Neumann algebra $M$, ruling out nonzero (even non-unital) ring homomorphisms from $S$ into $\U(M)$ or characteristic-zero matrix rank ultraproducts, and unital ring homomorphisms into any unital $C^*$-algebra admitting a tracial state (Theorem~\ref{thm:vn-main}).
\end{abstract}

\section{Introduction and main results}
\label{sec:intro}

\subsection{Background and Elek's question}

All rings in this paper are associative and unital, and a ring homomorphism is unital unless explicitly stated otherwise. Recall that a ring $S$ is \emph{von Neumann regular} (briefly, \emph{regular}) if for every $a \in S$ there exists $b \in S$ with $aba = a$. Following Goodearl~\cite[Chapter~16]{Goodearl} and Elek~\cite{Elek2013}, a \emph{rank function} on a ring $S$ is a map $\rk\colon S \to [0,1]$ satisfying:
\begin{enumerate}[label=(R\arabic*),leftmargin=2.8em]
\item $\rk(1) = 1$, and $\rk(x) = 0$ if and only if $x = 0$ (faithfulness);
\item $\rk(x + y) \le \rk(x) + \rk(y)$ for all $x, y \in S$;
\item $\rk(xy) \le \min\{\rk(x), \rk(y)\}$ for all $x, y \in S$;
\item $\rk(e + f) = \rk(e) + \rk(f)$ for all orthogonal idempotents $e, f \in S$ (meaning $e^2 = e$, $f^2 = f$, $ef = fe = 0$).
\end{enumerate}
A \emph{regular rank ring} is a von Neumann regular ring equipped with a rank function; a \emph{regular rank $K$-algebra} is a regular rank ring that is also an algebra over a field $K$. On a matrix ring $M_N(F)$ over a field $F$, the normalized matrix rank $A \mapsto \rank(A)/N$ is a rank function.

Given an index set $I$, an ultrafilter $\omega$ on $I$, positive integers $d_i$ ($i \in I$), and fields $F_i$ ($i \in I$), consider the product ring $\prod_{i \in I} M_{d_i}(F_i)$. The subset
\[
J_\omega := \Bigl\{(a_i)_{i \in I} \in \prod_{i \in I} M_{d_i}(F_i) : \lim_{i \to \omega} \frac{\rank(a_i)}{d_i} = 0\Bigr\}
\]
is a two-sided ideal, by subadditivity and multiplicativity of matrix rank. The quotient
\[
\RUF{\omega}{d_i}{F_i} := \Bigl(\prod_{i \in I} M_{d_i}(F_i)\Bigr) / J_\omega
\]
is the \emph{rank ultraproduct} of the matrix rings $M_{d_i}(F_i)$ along $\omega$~\cite{ElekSzabo,Elek2013}. It is itself a regular rank ring with rank function $\rk([(a_i)]) = \lim_{i \to \omega} \rank(a_i)/d_i$.

In~\cite[Section~1.2]{Elek2013}, Elek takes $I = \N$, a nonprincipal ultrafilter $\nu$ on $\N$, strictly increasing dimensions $d_1 < d_2 < \dots$, and $F_i = \C$ for all $i$, denoting the resulting rank ultraproduct by $M^{\mathrm{alg}}_\nu$. Calling the problem ``a purely algebraic version of the Connes Embedding Problem,'' Elek asks verbatim: ``One can ask of course, whether any countable dimensional regular rank ring embeds to some $M^{\mathrm{alg}}_\nu$.''

In the terminology of Arzhantseva and P\u{a}unescu~\cite[Definition~4.2]{AP2017}, a countably generated unital algebra $A$ over a field $F$ is \emph{linear sofic} over $F$ if it admits an injective unital $F$-algebra homomorphism into a rank ultraproduct $\prod_\omega^{\mathrm{rk}} M_{d_k}(F)$, while by~\cite[Proposition~11.2]{AP2017} the existence of a (not necessarily injective) unital $F$-algebra homomorphism into such an ultraproduct characterizes having almost finite-dimensional representations. Excluding every unital ring homomorphism into $\prod_\omega^{\mathrm{rk}} M_{d_i}(F_i)$ rules out both properties at once.

\subsection{Main results}

We distinguish explicitly between two tiers of results in this paper: the unital non-embedding results in Theorems~\ref{thm:main-ultra} and~\ref{thm:finite-form} (Sections~\ref{sec:ineq}--\ref{sec:proofs}) are formalized and machine-checked in Lean~4 (with the exact relation between the standalone specification file and the internal theorems detailed below and in Section~\ref{sec:formal}); Theorem~\ref{thm:vn-main} (including part~(c) on non-unital homomorphisms into matrix rank ultraproducts) and Proposition~\ref{prop:stably-finite} (Sections~\ref{sec:intro} and~\ref{sec:vn}) are proved on paper and have not been formalized in Lean~4.

\begin{theorem}[Lean~4: \code{elek\_countable}, \code{elek\_complex}]
\label{thm:main-ultra}
\leavevmode
\begin{enumerate}[label=\textup{(\alph*)},leftmargin=2.2em]
\item There exists a countable von Neumann regular $\Q$-algebra $S$ equipped with a faithful rank function $\rk\colon S \to [0,1]$ such that for every index set $I$, every ultrafilter $\omega$ on $I$, every family of positive integers $(d_i)_{i \in I}$ ($d_i \ge 1$ for all $i \in I$), and every family of fields $(F_i)_{i \in I}$ of characteristic zero, there is no unital ring homomorphism $S \to \RUF{\omega}{d_i}{F_i}$.
\item There exists a von Neumann regular $\C$-algebra $R$ equipped with a faithful rank function $\rk\colon R \to [0,1]$ with the same property: for every $I$, $\omega$, positive integers $(d_i)_{i \in I}$ ($d_i \ge 1$), and characteristic-zero fields $(F_i)_{i \in I}$, there is no unital ring homomorphism $R \to \RUF{\omega}{d_i}{F_i}$.
\end{enumerate}
\end{theorem}

In particular, neither $S$ nor $R$ admits any unital ring homomorphism into Elek's ultraproducts $M^{\mathrm{alg}}_\nu$. Note that the hypothesis that $S$ be a $\Q$-algebra is essential in part~(a): without it, the finite field $\F_2$ is a countable regular rank ring admitting no unital ring homomorphism into any ring containing $\Q$. In part~(b), the $\C$-algebra structure is obtained via a non-constructive Steinitz isomorphism $\overline{\Q}_7 \cong \C$ (depending on the Axiom of Choice; Lean: \code{nonempty\_padicAlgCl\_equiv\_complex}), and no claim is made regarding the dimension of $R$ over $\C$ (see Section~\ref{sec:open-readings}).

Our obstruction arises from an explicit finite presentation over $\Z$ built from the finite projective plane over $\F_\ell$. Let $\ell$ be a prime, let $(\PP, \LL)$ be the projective plane $\mathbb{P}^2(\F_\ell)$ with $n = \ell^2 + \ell + 1$ points and lines, and fix a bijection $\sigma\colon \PP \xrightarrow{\sim} \LL$ without fixed incidence, meaning $p \notin \sigma(p)$ for every $p \in \PP$ (such a $\sigma$ exists by Hall's marriage theorem; Lemma~\ref{lem:hall-sigma}). Set
\[
E := \{(p, p') \in \PP \times \PP : p \in \sigma(p')\}, \qquad |E| = (\ell + 1)n,
\]
and let $\mathrm{Gn}_\sigma := \{f_p : p \in \PP\} \cup \{x_{e,i}, y_{e,i} : e \in E,\ i \in \{0,1,2\}\}$ be a finite generator set. Define the finitely presented ring $B_\ell^\sigma$ by the generators $\mathrm{Gn}_\sigma$ subject to the finite set of integer polynomial relations
\begin{equation}
\label{eq:B-ell}
f_p^2 = f_p, \qquad f_p x_{e,i} = x_{e,i}, \qquad f_{p'} x_{e,i} = 0, \qquad \sum_{i=0}^{2} y_{e,i} x_{e,i} = 1
\end{equation}
for all $p \in \PP$, all $e = (p, p') \in E$, and all $i \in \{0,1,2\}$. (The condition $p \notin \sigma(p)$ is necessary: if $(p,p) \in E$, then~\eqref{eq:B-ell} forces $x_{(p,p),i} = f_p x_{(p,p),i} = 0$ and hence $1 = 0$.)

\begin{theorem}[Lean~4: checked \code{finite\_form}; internal \code{finite\_form\_paper}, \code{eps\_char0\_paper}, \code{exists\_config\_modL}, \code{relP\_of\_config}]
\label{thm:finite-form}
For $\ell = 7$ (so $n = 57$ and $|E| = 456$) and every fixed-incidence-free bijection $\sigma$ of $\mathbb{P}^2(\F_7)$, let $P_1, \dots, P_m \in \Z\langle \mathrm{Gn}_\sigma \rangle$ be the noncommutative integer polynomials corresponding to the relations~\eqref{eq:B-ell} of $B_7^\sigma$. Then there exists a constant $\varepsilon > 0$ such that:
\begin{enumerate}[label=\textup{(\alph*)},leftmargin=2.2em]
\item The relations $P_j = 0$ hold exactly in a countable regular rank $\Q$-algebra $S_\sigma$ satisfying Theorem~\ref{thm:main-ultra}\textup{(a)} (and, for the bijection $\sigma$ fixed in Theorem~\ref{thm:main-ultra}\textup{(a)}, in $S = S_\sigma$ itself).
\item For every field $F$ of characteristic zero, every integer $N \ge 1$, and every choice of matrices $A\colon \mathrm{Gn}_\sigma \to M_N(F)$, there exists an index $j$ with
\[
\rank\bigl(P_j(A)\bigr) > \varepsilon N.
\]
\item By contrast, for every integer $k \ge 0$, the exact same relations $P_j = 0$ admit an exact solution in $M_{3^{3k+1}}(\F_7)$. (More generally, for every prime $\ell$, every fixed-incidence-free $\sigma$, and every $k \ge 0$, the relations of $B_\ell^\sigma$ admit an exact solution in $M_{3^{3k+1}}(\F_\ell)$.)
\end{enumerate}
\end{theorem}

In Lean~4, \code{Challenge/Statement.lean} states \code{finite\_form} in existential form over $\C$, verified identical to \code{Solution.lean} by \code{Check/SameStatement.lean}: it asserts the existence of a finite presentation \code{(Gn, Rl, rel)}, a constant $\varepsilon > 0$, and a countable regular rank $\Q$-algebra $S$ in which the relations hold exactly, while no complex $N \times N$ matrices ($N \ge 1$) satisfy all relations up to rank $\varepsilon N$. The stronger forms stated in Theorem~\ref{thm:finite-form} are proved as internal theorems under \code{RankRings/} (built and axiom-checked, though not compared by \code{Check/SameStatement.lean}), quantifying over fields in Lean's universe \code{Type}: in \code{FiniteFormF.lean}, \code{finite\_form\_paper} proves parts~(a)--(b) for the literal relations~\eqref{eq:B-ell} of $B_7^\sigma$ for an existentially fixed $\sigma$ uniformly over all characteristic-zero fields $F$, while \code{eps\_char0\_paper} proves part~(b) for every fixed-incidence-free $\sigma$ and \code{finite\_form\_char0} gives the existential all-characteristic-zero-fields form; in \code{ModL.lean}, \code{exists\_config\_modL} and \code{rel\_modL} construct exact solutions of the equivalent presentation $\widetilde{B}_\ell^\sigma$ (\code{rel $\sigma$}) in $M_{3^{3k+1}}(\F_\ell)$ for every prime $\ell$, every fixed-incidence-free $\sigma$, and every $k \ge 0$ (yielding exact solutions of the literal relations $B_\ell^\sigma$ of part~(c) by composing with Lemma~\ref{lem:B-equiv}(a), Lean: \code{relP\_of\_config}), and \code{char\_contrast} pairs those characteristic-$7$ solutions of $\widetilde{B}_7^\sigma$ with the complex-matrix $\varepsilon$-gap. Note that quantifying over fields in universe $0$ (\code{Type}) in Theorems~\ref{thm:main-ultra} and~\ref{thm:finite-form} is mathematically harmless: any homomorphism out of a countable ring $S$, or any tuple of $N \times N$ matrices for a finite presentation, involves at most countably many matrix entries and therefore factors through countable subfields $F'_i \subseteq F_i$ (which lie in universe $0$), and matrix rank is invariant under field extension.

Extending the matrix incidence inequality of Section~\ref{sec:ineq} to finite von Neumann algebras yields a parallel operator-algebraic obstruction. Because Mathlib does not currently contain traces on von Neumann algebras or rings of affiliated operators, the following theorem (including part~(c) on non-unital maps into matrix rank ultraproducts) is proved on paper in Section~\ref{sec:vn} and is not checked in Lean~4.

\begin{theorem}[Paper proof, not formalized in Lean~4]
\label{thm:vn-main}
Let $A$ be the countable regular rank $\Q$-algebra $S$ of Theorem~\ref{thm:main-ultra}\textup{(a)}, or more generally any unital ring containing a solution of~\eqref{eq:B-ell} for $\ell = 7$ and some fixed-incidence-free $\sigma$.
\begin{enumerate}[label=\textup{(\alph*)},leftmargin=2.2em]
\item For every nonzero finite von Neumann algebra $M$, there is no nonzero (even non-unital) ring homomorphism $A \to \U(M)$, where $\U(M)$ is the $*$-regular ring of closed, densely defined operators affiliated with $M$.
\item For every unital $C^*$-algebra $C$ admitting a tracial state, there is no unital ring homomorphism $A \to C$.
\item For every index set $I$, ultrafilter $\omega$ on $I$, positive integers $(d_i)_{i \in I}$ ($d_i \ge 1$), and fields $(F_i)_{i \in I}$ of characteristic zero, there is no nonzero (even non-unital) ring homomorphism $A \to \RUF{\omega}{d_i}{F_i}$.
\end{enumerate}
\end{theorem}

Neither Theorem~\ref{thm:main-ultra} nor Theorem~\ref{thm:vn-main}(a) is derived from the other here: a rank ultraproduct of matrix rings is not presented as a ring $\U(M)$, and $\U(M)$ is not assumed to embed into a rank ultraproduct of matrix rings.

\subsection{Stable finiteness, Sylvester rank functions, and contrast with Greenfeld}

Recall that a ring $A$ is \emph{directly finite} if $xy = 1$ implies $yx = 1$ for all $x, y \in A$, and \emph{stably finite} if $M_d(A)$ is directly finite for every $d \ge 1$. Following Schofield~\cite[Chapter~7]{Schofield1985}, a \emph{Sylvester matrix rank function} on a ring $A$ is a map $\rk$ assigning to every rectangular matrix $X$ over $A$ a real number $\rk(X) \ge 0$ such that: (i)~$\rk(0) = 0$ and $\rk(1) = 1$; (ii)~$\rk(XY) \le \min\{\rk(X), \rk(Y)\}$ for conformable matrices; (iii)~$\rk\bigl(\begin{smallmatrix} X & 0 \\ 0 & Y \end{smallmatrix}\bigr) = \rk(X) + \rk(Y)$; and (iv)~$\rk\bigl(\begin{smallmatrix} X & Z \\ 0 & Y \end{smallmatrix}\bigr) \ge \rk(X) + \rk(Y)$ for blocks of matching sizes. It is \emph{faithful} if $\rk(x) > 0$ for every nonzero $x \in A$.

Resolving a question of Arzhantseva and P\u{a}unescu~\cite{AP2017}, Greenfeld~\cite{Greenfeld2023} constructed finitely generated stably finite algebras over every field that are not linear sofic. Greenfeld's primary obstruction criterion~\cite[Lemma~2.1]{Greenfeld2023} certifies non-linear-soficity by exhibiting nonzero elements $x, y, z \in A$ satisfying $x \in yxA$, $z \in xA$, and $yz = 0$; this criterion is used for the stably finite algebras constructed in~\cite{Greenfeld2023}. As the next proposition shows, this mechanism is incompatible with having a faithful Sylvester matrix rank function, whereas our algebras possess one.

\begin{proposition}[Paper argument]
\label{prop:stably-finite}
\leavevmode
\begin{enumerate}[label=\textup{(\alph*)},leftmargin=2.2em]
\item The countable regular rank $\Q$-algebra $S$ of Theorem~\ref{thm:main-ultra}\textup{(a)} is stably finite, admits a faithful Sylvester matrix rank function extending $\rk$, and is not linear sofic over $\Q$. Likewise, the $\C$-subalgebra $A \subseteq R$ of Theorem~\ref{thm:main-ultra}\textup{(b)} generated by the solution of~\eqref{eq:B-ell} is a finitely generated stably finite $\C$-algebra admitting a faithful Sylvester matrix rank function that is not linear sofic over $\C$.
\item Conversely, if an algebra $A$ contains nonzero elements $x, y, z \in A$ with $x \in yxA$, $z \in xA$, and $yz = 0$ (the criterion of~\cite[Lemma~2.1]{Greenfeld2023}), then every Sylvester matrix rank function on $A$ satisfies $\rk(z) = 0$; in particular, $A$ admits no faithful Sylvester matrix rank function and admits no unital embedding into any regular rank ring.
\end{enumerate}
\end{proposition}

\begin{proof}
(a) On any von Neumann regular ring $T$ with a faithful rank function $\rk$, every finitely generated right submodule of $T^d$ is projective and a direct summand of $T^d$~\cite[Theorem~1.11]{Goodearl}, and $\rk$ extends uniquely to a faithful Sylvester matrix rank function on $T$ (equivalently, a faithful rank function $\rk_d$ on $M_d(T)$ for every $d \ge 1$; see~\cite[Proposition~16.1, Corollary~16.10]{Goodearl} and~\cite[Chapter~7]{Schofield1985}). In a ring with a faithful rank function $\rk_d$, if $xy = 1$ then $e = yx$ is an idempotent satisfying $1 = \rk_d(1) = \rk_d(xy) = \rk_d(x(yx)y) \le \rk_d(e) \le 1$, so $\rk_d(1 - e) = \rk_d(1) - \rk_d(e) = 0$ and hence $yx = e = 1$. Thus $T$ (and every unital subring of $T$) is stably finite and inherits a faithful Sylvester matrix rank function.

For non-linear-soficity, observe directly that both $S$ and the finitely generated $\C$-subalgebra $A \subseteq R$ contain an exact solution of~\eqref{eq:B-ell} for $\ell = 7$. For any field $F$ of characteristic zero (in particular $F = \Q$ for $S$, and $F = \C$ for $A$), any rank ultraproduct $Q = \RUF{\omega}{d_k}{F}$ ($d_k \ge 1$) is von Neumann regular (Section~\ref{sec:proofs}). A unital ring homomorphism from $S$ or $A$ into $Q$ would carry the solution of~\eqref{eq:B-ell} into $Q$; by Lemma~\ref{lem:B-equiv}(b) it would yield a solution of~\eqref{eq:B-tilde} in $Q$, contradicting Proposition~\ref{prop:no-config}. Hence $S$ is not linear sofic over $\Q$ (and admits no unital map into $\RUF{\omega}{d_k}{\Q}$), and $A$ is not linear sofic over $\C$ (and admits no unital map into $\RUF{\omega}{d_k}{\C}$).

(b) Write $x = yxa$ and $z = xb$ with $a, b \in A$. Since $yz = 0$, we have the matrix factorization
\[
\begin{pmatrix} yx & 0 \\ x & z \end{pmatrix}
= \begin{pmatrix} y \\ 1 \end{pmatrix} x \begin{pmatrix} 1 & b \end{pmatrix}.
\]
Swapping rows and columns via invertible permutation matrices and applying axioms~(ii) and~(iv) of a Sylvester matrix rank function gives
\[
\rk(z) + \rk(yx)
\le \rk\begin{pmatrix} z & x \\ 0 & yx \end{pmatrix}
= \rk\begin{pmatrix} yx & 0 \\ x & z \end{pmatrix}
\le \rk(x).
\]
On the other hand, $x = yxa$ gives $\rk(x) \le \rk(yx)$ by axiom~(ii). Therefore $\rk(z) + \rk(x) \le \rk(x)$, so $\rk(z) = 0$. Since $z \neq 0$, no Sylvester matrix rank function on $A$ can be faithful, and since any unital embedding into a regular rank ring pulls back a faithful Sylvester matrix rank function, $A$ admits no unital embedding into any regular rank ring.
\end{proof}

\subsection{Readings left open and limitations}
\label{sec:open-readings}

Theorem~\ref{thm:main-ultra} answers Elek's question in the reading ``countable regular rank $\Q$-algebra with unital homomorphisms'' (ruling out non-unital embeddings as well relies on the paper-only Theorem~\ref{thm:vn-main}(c)). Three other natural readings remain open:
\begin{enumerate}[label=\textup{(\arabic*)},leftmargin=2.2em]
\item \emph{Regular rank $\C$-algebras of countable dimension over $\C$.} Every element $a$ of a countable-dimensional regular rank $\C$-algebra $A$ is algebraic over $\C$. Indeed, if $a - \lambda$ ($\lambda \in \C$) is not invertible, choose $b \in A$ with $(a - \lambda)b(a - \lambda) = a - \lambda$; since $A$ is directly finite (as in Proposition~\ref{prop:stably-finite}(a)), $b(a - \lambda) \neq 1$, so $r_\lambda := 1 - b(a - \lambda)$ is a nonzero element satisfying $a r_\lambda = \lambda r_\lambda$. Eigenvectors of left multiplication by $a$ for distinct eigenvalues are $\C$-linearly independent, so $a - \lambda$ is invertible for all but countably many $\lambda \in \C$. Since $\dim_\C A \le \aleph_0$, the uncountably many inverses $(a - \lambda)^{-1}$ satisfy a nontrivial linear relation $\sum_{j=1}^m c_j (a - \lambda_j)^{-1} = 0$ with distinct $\lambda_j$, and multiplying by $\prod_{j=1}^m (a - \lambda_j)$ yields a nonzero polynomial $p \in \C[t]$ (satisfying $p(\lambda_k) = c_k \prod_{j \neq k}(\lambda_k - \lambda_j) \neq 0$ whenever $c_k \neq 0$) with $p(a) = 0$. Our construction in Theorem~\ref{thm:main-ultra}(b) uses $\overline{\Q}_7 \cong \C$ and has uncountable $\C$-dimension; it does not decide whether a countable-dimensional regular rank $\C$-algebra can fail to embed into $M^{\mathrm{alg}}_\nu$.
\item \emph{$*$-regular rings with involution-preserving homomorphisms.} Elek's ultraproducts $M^{\mathrm{alg}}_\nu$ carry a natural involution, whereas our rings $S$ and $R$ carry no involution in Lean~4; the $*$-regular formulation with $*$-homomorphisms is not treated here.
\item \emph{Rank ultraproducts over fields of positive characteristic.} Our characteristic-zero obstruction does not rule out unital maps from $S$ into rank ultraproducts $\RUF{\omega}{d_i}{\F_{p_i}}$ with $p_i \to \infty$ (which are also $\Q$-algebras).
\end{enumerate}
Finally, our results have no direct implications for groups, the Kaplansky conjectures, or the Connes Embedding Problem: for every group $G$, the group algebra $\C[G]$ embeds into $\U(L(G))$, where the incidence inequality of Theorem~\ref{thm:vn-ineq} holds, so no matrix ring over $\C[G]$ contains a solution of $B_7^\sigma$.

\subsection{Overview of the proof}

The argument combines five ingredients. First (Section~\ref{sec:ineq}), for any finite projective plane $(\PP,\LL)$ of order $q$ with $n = q^2+q+1$ points, any bijection $\sigma\colon \PP \xrightarrow{\sim} \LL$, and any idempotents $f_p \in M_N(\C)$, we prove the average incidence rank bound
\[
\frac{1}{|E|}\sum_{(p,p') \in E} \frac{\dim(\im f_p \cap \ker f_{p'})}{N} \le \frac{1 + \sqrt{q}/(q+1)}{4}
\]
by balancing the inner product via an elementary algebraic log-determinant descent (related to the scaling methods of Barthe~\cite{Barthe1998}, Barak--Dvir--Wigderson--Yehudayoff~\cite{BDWY2011}, and Dvir--Hu~\cite{DvirHu2016}) and applying expander mixing on the projective plane. Second (Section~\ref{sec:ultrametric}), over any ultrametric normed field $K$, we define the scale-$r$ rank $\srk_r(A)$ and use Smith normal form over $\OO_K$ to show that tight matrix sequences modulo null sequences form a von Neumann regular rank $K$-algebra $\RR(K, \iota, \omega)$. Third (Section~\ref{sec:exotic}), $(3k+1)$-fold tensor powers of $\ell$-adic lifts of adapted flags in $\mathbb{P}^2(\F_\ell)$ solve $B_\ell^\sigma$ exactly in $\RR(K, \iota, \omega)$ (forcing every edge intersection rank to be at least $1/3$), while their reductions modulo $\ell$ solve $B_\ell^\sigma$ exactly in $M_{3^{3k+1}}(\F_\ell)$. Fourth (Section~\ref{sec:proofs}), at $\ell = 7$ we have $(1 + \sqrt{7}/8)/4 = (8+\sqrt{7})/32 \approx 0.33268 < 1/3$ (equivalently $63 < 64$); lifting approximate solutions in any characteristic-zero matrix ultraproduct and embedding a countable subfield into $\C$ gives a contradiction, yielding Theorems~\ref{thm:main-ultra} and~\ref{thm:finite-form}. Finally (Section~\ref{sec:vn}), replacing the log-determinant potential by $\sum_p \tau(\log L(g f_p g^{-1}))$ and truncating unbounded idempotents extends the incidence bound to $\U(M)$, giving Theorem~\ref{thm:vn-main}.

\section{Projective plane incidence rank bound in characteristic zero}
\label{sec:ineq}

In this section we prove the finite-dimensional incidence rank bound over $\C$ (Theorem~\ref{thm:ineq-C}); all results correspond to the Lean modules \code{LogDet.lean}, \code{Descent.lean}, \code{Plane.lean}, and \code{Ineq.lean} in \code{RankRings/}.

\subsection{Log-determinant inequalities and algebraic descent}

Fix $N \ge 1$. For $X \in M_N(\C)$, let $X^*$ denote its conjugate transpose and define $\ld(X) := \log |\det X| \in \R$ whenever $\det X \neq 0$ (Lean: \code{RankRings.ld}). Clearly $\ld(XY) = \ld(X) + \ld(Y)$ for invertible $X, Y$, and $\ld(X^{-1}) = -\ld(X)$. We write $X \succ 0$ (resp.\ $X \succeq 0$) to indicate that $X$ is Hermitian positive definite (resp.\ positive semidefinite).

\begin{lemma}[Lean: \code{ld\_le\_trace}, \code{ld\_add\_ld\_le}]
\label{lem:ld-trace}
For all $X, W \in M_N(\C)$ with $X \succ 0$ and $W \succ 0$,
\[
\ld(X) \le \RePart\tr(X) - N
\qquad\text{and}\qquad
\ld(X) + \ld(W) \le \RePart\tr(XW) - N.
\]
\end{lemma}

\begin{proof}
Let $\lambda_1, \dots, \lambda_N > 0$ be the eigenvalues of $X$. Since $\log t \le t - 1$ for all $t > 0$,
\[
\ld(X) = \sum_{j=1}^N \log \lambda_j \le \sum_{j=1}^N (\lambda_j - 1) = \RePart\tr(X) - N.
\]
For the second inequality, write $W = B^* B$ with $B \in \GL_N(\C)$. Then $B X B^* \succ 0$, $\ld(B X B^*) = \ld(X) + \ld(B^* B) = \ld(X) + \ld(W)$, and $\tr(B X B^*) = \tr(X B^* B) = \tr(XW)$. Applying the first inequality to $B X B^*$ yields the result.
\end{proof}

\begin{lemma}[Lean: \code{scalar\_step}, \code{step\_ineq}]
\label{lem:step-ineq}
\leavevmode
\begin{enumerate}[label=\textup{(\alph*)},leftmargin=2.2em]
\item For every real $x$ with $|x| \le \frac{1}{4}$,
\[
(1 - x)^2 (4x + 1) - 1 - 2\log(1 - x) \le 4x - 2x^2.
\]
\item Let $G \in M_N(\C)$ be Hermitian with $\tr(G) = 0$, and let $s > 0$ satisfy $-sI \preceq G \preceq sI$. Put $T := I - \frac{1}{4s} G$. Then $T$ is Hermitian, $\det T \neq 0$, and
\[
\RePart\tr\bigl(T^2 (G + sI)\bigr) - sN - 2s\ld(T) \le -\frac{\RePart\tr(G^2)}{8s}.
\]
\end{enumerate}
\end{lemma}

\begin{proof}
(a) Since $1 - x \ge \frac{3}{4} > 0$, the inequality $\log t \le t - 1$ at $t = (1 - x)^{-1}$ gives $-\log(1 - x) \le (1 - x)^{-1} - 1$. Moreover,
\[
(1 - x)\Bigl(1 + x + \frac{4}{3}x^2\Bigr) = 1 + \frac{4}{3}x^2\Bigl(\frac{1}{4} - x\Bigr) \ge 1,
\]
so $-\log(1 - x) \le x + \frac{4}{3}x^2$. Expanding $(1 - x)^2(4x + 1) - 1 = 2x - 7x^2 + 4x^3$, we obtain
\[
(1 - x)^2(4x + 1) - 1 - 2\log(1 - x)
\le 4x - \frac{13}{3}x^2 + 4x^3
= 4x - 2x^2 - x^2\Bigl(\frac{7}{3} - 4x\Bigr)
\le 4x - 2x^2,
\]
as $4x \le 1 < \frac{7}{3}$.

(b) Diagonalize $G = U \diag(\gamma_1, \dots, \gamma_N) U^*$ with $U$ unitary and $\gamma_j \in [-s, s]$. Set $x_j := \gamma_j / (4s)$, so $|x_j| \le \frac{1}{4}$, $\gamma_j = 4s x_j$, and $T = U \diag(1 - x_1, \dots, 1 - x_N) U^*$. Thus $T$ is Hermitian with eigenvalues $1 - x_j \ge \frac{3}{4} > 0$, so $\det T \neq 0$ and $\ld(T) = \sum_{j=1}^N \log(1 - x_j)$. By~(a) multiplied by $s > 0$,
\[
(1 - x_j)^2(\gamma_j + s) - s - 2s\log(1 - x_j)
= s\bigl((1 - x_j)^2(4x_j + 1) - 1 - 2\log(1 - x_j)\bigr)
\le s(4x_j - 2x_j^2) = \gamma_j - \frac{\gamma_j^2}{8s}.
\]
Summing over $j = 1, \dots, N$ and using $\sum_{j=1}^N \gamma_j = \tr(G) = 0$ and $\sum_{j=1}^N \gamma_j^2 = \RePart\tr(G^2)$ gives the claimed inequality.
\end{proof}

\subsection{The canonical metric of an idempotent and balancing}

For any matrix $E \in M_N(\C)$, define (Lean: \code{RankRings.LL}, \code{RankRings.pr})
\[
L(E) := E^* E + (I - E)^*(I - E), \qquad \pr(E) := E L(E)^{-1}.
\]

\begin{lemma}[Lean: \code{Descent.lean}, \code{comparison}, \code{comparison\_step}]
\label{lem:L-idem}
Let $E \in M_N(\C)$ be an idempotent ($E^2 = E$).
\begin{enumerate}[label=\textup{(\alph*)},leftmargin=2.2em]
\item $L(E) = L(I - E) \succ 0$, and $L(E)E = E^*E = E^* L(E)$.
\item $\pr(E) = L(E)^{-1} E^*$ is the orthogonal projection onto $\im(E)$: specifically, $\pr(E)^* = \pr(E)$, $\pr(E)^2 = \pr(E)$, $\pr(E)E = E$, $E\pr(E) = \pr(E)$, and $\tr(\pr(E)) = \tr(E)$. Similarly, $\pr(I - E)$ is the orthogonal projection onto $\ker(E)$.
\item $\pr(E) + \pr(I - E) = L(E)^{-1}$, $\tr(L(E)^{-1}) = N$, and $\ld(L(E)) \ge 0$.
\item If $Z \in M_N(\C)$ satisfies $Z \succ 0$ and $EZ = ZE^*$, then
\[
\ld(L(E)) \le \RePart\tr(Z) - N - \ld(Z).
\]
\item For every invertible Hermitian $T \in M_N(\C)$,
\[
\ld\bigl(L(T E T^{-1})\bigr) \le \ld(L(E)) + \RePart\tr\bigl(T^2 L(E)^{-1}\bigr) - N - 2\ld(T).
\]
\end{enumerate}
\end{lemma}

\begin{proof}
(a) For every nonzero $v \in \C^N$, $v^* L(E) v = \|Ev\|^2 + \|(I - E)v\|^2 > 0$ since $Ev + (I - E)v = v \neq 0$, so $L(E) \succ 0$. Clearly $L(I - E) = L(E)$. Since $(I - E)E = 0$, we have $L(E)E = E^*E^2 + (I-E)^*(I-E)E = E^*E$; taking conjugate transposes gives $E^* L(E) = E^* E$.

(b) By~(a), $E L(E)^{-1} = L(E)^{-1}(L(E)E)L(E)^{-1} = L(E)^{-1}(E^* L(E))L(E)^{-1} = L(E)^{-1}E^*$, so $\pr(E)^* = \pr(E)$. Hence $\pr(E)E = L(E)^{-1}E^*E = L(E)^{-1}L(E)E = E$, $E\pr(E) = E^2 L(E)^{-1} = \pr(E)$, and $\pr(E)^2 = \pr(E) E L(E)^{-1} = E L(E)^{-1} = \pr(E)$. Also $\tr(\pr(E)) = \tr(E\pr(E)) = \tr(\pr(E)E) = \tr(E)$. The same holds for the idempotent $I - E$, whose image is $\ker(E)$.

(c) Adding the definitions gives $\pr(E) + \pr(I - E) = (E + I - E)L(E)^{-1} = L(E)^{-1}$. Taking traces gives $\tr(L(E)^{-1}) = \tr(E) + \tr(I - E) = N$. By Lemma~\ref{lem:ld-trace} applied to $L(E)^{-1} \succ 0$,
\[
-\ld(L(E)) = \ld(L(E)^{-1}) \le \RePart\tr(L(E)^{-1}) - N = 0,
\]
so $\ld(L(E)) \ge 0$.

(d) Expanding $L(E) = I - E - E^* + 2E^*E$ and using $EZ = ZE^*$ together with cyclicity of trace gives $\tr(ZE^*) = \tr(EZ) = \tr(ZE)$ and
\[
\tr(Z E^* E) = \tr(E Z E) = \tr(E^2 Z) = \tr(EZ) = \tr(ZE).
\]
Thus $\tr(Z L(E)) = \tr(Z) - \tr(ZE) - \tr(ZE^*) + 2\tr(ZE^*E) = \tr(Z)$. Lemma~\ref{lem:ld-trace} applied to $Z$ and $L(E)$ yields $\ld(Z) + \ld(L(E)) \le \RePart\tr(Z) - N$.

(e) Put $E' := T E T^{-1}$ (which is idempotent) and $Z := T L(E)^{-1} T \succ 0$. Since $T^* = T$ and $(T^{-1})^* = T^{-1}$, we verify
\[
E' Z = T E L(E)^{-1} T = T L(E)^{-1} E^* T = (T L(E)^{-1} T)(T^{-1} E^* T) = Z (E')^*.
\]
Applying~(d) to $E'$ and $Z$, using $\tr(Z) = \tr(T^2 L(E)^{-1})$ and $\ld(Z) = 2\ld(T) - \ld(L(E))$, yields~(e).
\end{proof}

\begin{proposition}[Almost-balanced inner product; Lean: \code{RankRings.exists\_balanced}]
\label{prop:balanced}
Let $\PP$ be a nonempty finite set, let $f_p \in M_N(\C)$ ($p \in \PP$) be idempotents, and let $\eta > 0$. Then there exists $g \in \GL_N(\C)$ such that, writing $E_p := g f_p g^{-1}$ and
\[
G := \sum_{p \in \PP} \bigl(L(E_p)^{-1} - I\bigr),
\]
we have $G^* = G$ and $\RePart\tr(G^2) \le \eta$.
\end{proposition}

\begin{proof}
Set $s := |\PP| > 0$, and consider the potential $\Phi(g) := \sum_{p \in \PP} \ld(L(g f_p g^{-1}))$ for $g \in \GL_N(\C)$. By Lemma~\ref{lem:L-idem}(c), $\Phi(g) \ge 0$ for all $g \in \GL_N(\C)$, so $\inf \Phi \ge 0$ exists. Choose $g \in \GL_N(\C)$ such that
\[
\Phi(g) < \inf_{h \in \GL_N(\C)} \Phi(h) + \frac{\eta}{8s}.
\]
With $E_p = g f_p g^{-1}$ and $G = \sum_{p \in \PP}(L(E_p)^{-1} - I)$, we verify the hypotheses of Lemma~\ref{lem:step-ineq}(b):
\begin{itemize}[leftmargin=1.8em]
\item $G$ is Hermitian, and $\tr(G) = \sum_{p \in \PP}(\tr(L(E_p)^{-1}) - N) = 0$ by Lemma~\ref{lem:L-idem}(c).
\item $G + sI = \sum_{p \in \PP} L(E_p)^{-1} \succeq 0$.
\item $sI - G = \sum_{p \in \PP}\bigl(2I - L(E_p)^{-1}\bigr) = \sum_{p \in \PP}\bigl((I - \pr(E_p)) + (I - \pr(I - E_p))\bigr) \succeq 0$, since $I - \pr(E_p)$ and $I - \pr(I - E_p)$ are orthogonal projections.
\end{itemize}
Set $T := I - \frac{1}{4s}G$. By Lemma~\ref{lem:step-ineq}(b), $T$ is invertible and Hermitian. Applying Lemma~\ref{lem:L-idem}(e) to each $E_p$ and summing over $p \in \PP$ yields
\begin{align*}
\Phi(Tg)
&= \sum_{p \in \PP} \ld\bigl(L(T E_p T^{-1})\bigr) \\
&\le \Phi(g) + \RePart\tr\bigl(T^2 (G + sI)\bigr) - sN - 2s\ld(T)
\le \Phi(g) - \frac{\RePart\tr(G^2)}{8s}.
\end{align*}
Since $\Phi(Tg) \ge \inf \Phi > \Phi(g) - \frac{\eta}{8s}$, it follows that $\RePart\tr(G^2) \le \eta$.
\end{proof}

\subsection{Expander mixing on a projective plane and the incidence rank bound}

Let $(\PP, \LL)$ be a finite projective plane of order $q \ge 2$, with $n = |\PP| = |\LL| = q^2 + q + 1$. Each point lies on $q+1$ lines, each line contains $q+1$ points, and any two distinct points lie on a unique common line. Define the incidence indicator $\mathrm{inc}(p, l) = 1$ if $p \in l$ and $0$ otherwise. Then for all $p, p' \in \PP$,
\begin{equation}
\label{eq:inc-orth}
\sum_{l \in \LL} \mathrm{inc}(p, l)\,\mathrm{inc}(p', l) = 1 + q\,\delta_{p, p'}.
\end{equation}

\begin{lemma}[Hilbert-space mixing inequality; Lean: \code{RankRings.plane\_ineq}]
\label{lem:plane-ineq}
Let $(\PP, \LL)$ be a finite projective plane of order $q \ge 2$ with $n = q^2 + q + 1$ points, let $H$ be a real inner product space, and let $a\colon \PP \to H$ and $b\colon \LL \to H$ satisfy $\|a(p)\|^2 = \|b(l)\|^2 = \mu$ for all $p \in \PP$ and $l \in \LL$. Put $A := \sum_{p \in \PP} a(p)$ and $B := \sum_{l \in \LL} b(l)$. Then
\[
\sum_{l \in \LL} \sum_{p \in l} \langle a(p), b(l) \rangle
\le \sqrt{q}\, n \mu + \frac{q + 1}{n} \cdot \frac{\|A + B\|^2}{4}.
\]
\end{lemma}

\begin{proof}
For each $l \in \LL$, set $s(l) := \sum_{p \in \PP} \mathrm{inc}(p, l)\, a(p)$, put $c := (q + 1)/n$, and set $C(l) := s(l) - c A$. By~\eqref{eq:inc-orth},
\[
\sum_{l \in \LL} \|s(l)\|^2
= \sum_{p, p' \in \PP} (1 + q\,\delta_{p, p'})\langle a(p), a(p') \rangle
= q n \mu + \|A\|^2,
\]
and $\sum_{l \in \LL} s(l) = (q + 1) A$. Expanding $\|C(l)\|^2 = \|s(l)\|^2 - 2c\langle s(l), A\rangle + c^2\|A\|^2$ and summing over $l \in \LL$ gives
\[
\sum_{l \in \LL} \|C(l)\|^2
= q n \mu + \bigl(1 - 2c(q+1) + n c^2\bigr)\|A\|^2
= q n \mu + \Bigl(1 - \frac{(q+1)^2}{n}\Bigr)\|A\|^2
\le q n \mu,
\]
since $(q+1)^2 = q^2 + 2q + 1 > q^2 + q + 1 = n$. For each $l \in \LL$, the inequality $\|C(l) - \sqrt{q}\,b(l)\|^2 \ge 0$ rearranges to $2\sqrt{q}\langle C(l), b(l)\rangle \le \|C(l)\|^2 + q \mu$. Summing over $l \in \LL$ gives
\[
2\sqrt{q} \sum_{l \in \LL} \langle C(l), b(l) \rangle \le q n \mu + q n \mu = 2 q n \mu,
\]
so $\sum_{l \in \LL} \langle C(l), b(l) \rangle \le \sqrt{q}\, n \mu$. Furthermore, $\|A - B\|^2 \ge 0$ implies $\langle A, B \rangle \le \|A + B\|^2 / 4$. Combining these bounds in the identity
\[
\sum_{l \in \LL} \sum_{p \in l} \langle a(p), b(l) \rangle
= \sum_{l \in \LL} \langle C(l), b(l) \rangle + c\langle A, B \rangle
\]
completes the proof.
\end{proof}

\begin{theorem}[Incidence rank bound over $\C$; Lean: \code{RankRings.ineq\_C}]
\label{thm:ineq-C}
Let $(\PP, \LL)$ be a finite projective plane of order $q \ge 2$ with $n = q^2 + q + 1$ points, let $\sigma\colon \PP \xrightarrow{\sim} \LL$ be a bijection, and set $E := \{(p, p') \in \PP \times \PP : p \in \sigma(p')\}$. Let $f_p \in M_N(\C)$ ($p \in \PP$) be idempotents, and for each $(p, p') \in E$ let $r_{p, p'} \in M_N(\C)$ be an idempotent satisfying $f_p r_{p, p'} = r_{p, p'}$ and $f_{p'} r_{p, p'} = 0$. Then
\begin{equation}
\label{eq:ineq-C}
4 \sum_{(p, p') \in E} \RePart\tr(r_{p, p'}) \le (q + 1 + \sqrt{q})\, n N.
\end{equation}
In particular,
\begin{equation}
\label{eq:ineq-C-avg}
\frac{1}{|E|} \sum_{(p, p') \in E} \frac{\dim(\im f_p \cap \ker f_{p'})}{N} \le \frac{1 + \sqrt{q}/(q + 1)}{4}.
\end{equation}
\end{theorem}

\begin{proof}
Fix $\varepsilon > 0$ and apply Proposition~\ref{prop:balanced} with $\eta := \varepsilon n / (q + 1)$ to obtain $g \in \GL_N(\C)$ such that $E_p := g f_p g^{-1}$ and $G := \sum_{p \in \PP}(L(E_p)^{-1} - I)$ satisfy $G^* = G$ and $\RePart\tr(G^2) \le \eta$. Endow $M_N(\C)$ with the real Frobenius inner product $\langle X, Y \rangle := \RePart\tr(X^* Y)$, and define Hermitian involutions
\[
A_p := 2\pr(E_p) - I, \qquad B_{p'} := 2\pr(I - E_{p'}) - I \qquad (p, p' \in \PP).
\]
Since $A_p^2 = B_{p'}^2 = I$, we have $\|A_p\|^2 = \|B_{p'}\|^2 = N$. Moreover, by Lemma~\ref{lem:L-idem}(c),
\[
\sum_{p \in \PP} A_p + \sum_{l \in \LL} B_{\sigma^{-1}(l)}
= \sum_{p \in \PP} \bigl(2\pr(E_p) + 2\pr(I - E_p) - 2I\bigr)
= 2\sum_{p \in \PP} \bigl(L(E_p)^{-1} - I\bigr) = 2G,
\]
so $\|2G\|^2 / 4 = \RePart\tr(G^2) \le \varepsilon n / (q + 1)$. Applying Lemma~\ref{lem:plane-ineq} to $p \mapsto A_p$ and $l \mapsto B_{\sigma^{-1}(l)}$ yields
\begin{equation}
\label{eq:plane-applied}
\sum_{(p, p') \in E} \langle A_p, B_{p'} \rangle \le \sqrt{q}\, n N + \varepsilon.
\end{equation}
On the other hand, using $\tr(\pr(E_p)) = \tr(E_p) = \tr(f_p)$ and $\tr(\pr(I - E_{p'})) = \tr(I - f_{p'})$,
\[
\langle A_p, B_{p'} \rangle = 4\RePart\tr\bigl(\pr(E_p)\pr(I - E_{p'})\bigr) - 2\RePart\tr(f_p) - 2\RePart\tr(I - f_{p'}) + N.
\]
Summing over $(p, p') \in E$ (where every $p$ and every $p'$ appears $q+1$ times, and $|E| = (q+1)n$) and using $\sum_{p \in \PP}\bigl(\tr(f_p) + \tr(I - f_p)\bigr) = n N$ gives
\begin{equation}
\label{eq:main-sum}
\sum_{(p, p') \in E} \langle A_p, B_{p'} \rangle
= 4\sum_{(p, p') \in E} \RePart\tr\bigl(\pr(E_p)\pr(I - E_{p'})\bigr) - (q + 1) n N.
\end{equation}
Now fix $(p, p') \in E$ and put $X := g r_{p, p'} g^{-1}$. Then $X^2 = X$, $E_p X = X$, and $(I - E_{p'})X = X$. Setting $P := \pr(E_p)$, $Q := \pr(I - E_{p'})$, and $R := \pr(X)$, we have $P X = P E_p X = E_p X = X$ and thus $P R = P X L(X)^{-1} = X L(X)^{-1} = R$; similarly $Q R = R$. Since $Q - R = (Q - R)^*(Q - R) \succeq 0$, the matrix $P(Q - R)P \succeq 0$ has nonnegative real trace, so
\[
\RePart\tr(r_{p, p'}) = \RePart\tr(R) = \RePart\tr(PRP) \le \RePart\tr(PQP) = \RePart\tr(PQ)
\]
(Lean: \code{RankRings.trace\_mul\_ge}). Combining this with~\eqref{eq:plane-applied} and~\eqref{eq:main-sum} yields
\[
4\sum_{(p, p') \in E} \RePart\tr(r_{p, p'}) \le (q + 1 + \sqrt{q}) n N + \varepsilon.
\]
Since $\varepsilon > 0$ was arbitrary, \eqref{eq:ineq-C} follows. Taking $r_{p, p'}$ to be any projection onto $\im f_p \cap \ker f_{p'}$ (so $\tr(r_{p, p'}) = \dim(\im f_p \cap \ker f_{p'})$) and dividing by $4|E|N = 4(q+1)nN$ yields~\eqref{eq:ineq-C-avg}.
\end{proof}

\section{Ultrametric regular rank rings and the exotic configuration}
\label{sec:ultrametric}

In this section we construct the ultrametric regular rank ring $\RR(K, \iota, \omega)$ and build an exact solution of $B_\ell^\sigma$ inside it for every prime $\ell$ (Lean modules \code{SRank.lean}--\code{RankFn.lean}, \code{Words.lean}--\code{ExoticS.lean}, and \code{ModL.lean}; see Table~\ref{tab:lean-map}).

\subsection{Scale-\texorpdfstring{$r$}{r} rank and Smith normal form over an ultrametric field}

Let $K$ be a field equipped with a non-Archimedean (ultrametric) absolute value $\|\cdot\|$, meaning $\|x + y\| \le \max\{\|x\|, \|y\|\}$ for all $x, y \in K$. Let $\OO_K := \{x \in K : \|x\| \le 1\}$ be its valuation ring. A matrix $X \in M_{m \times n}(K)$ is \emph{bounded} if $X \in M_{m \times n}(\OO_K)$ (all entries have norm at most $1$); a square matrix $U \in M_n(K)$ is \emph{unimodular} if $U \in \GL_n(\OO_K)$ (both $U$ and $U^{-1}$ are bounded).

For any real $r \ge 0$ and any $A \in M_{m \times n}(K)$, define the \emph{rank at scale $r$} (Lean: \code{RankRings.srk}) by
\[
\srk_r(A) := \min\bigl\{\rank(A') : A' \in M_{m \times n}(K),\; \max_{i,j} \|A_{ij} - A'_{ij}\| \le r\bigr\}.
\]

\begin{lemma}[Lean: \code{srk\_add\_le}, \code{srk\_mul\_le}]
\label{lem:srk-basic}
For all conformable matrices $A, B$ over $K$ and all $r, s \ge 0$:
\begin{enumerate}[label=\textup{(\alph*)},leftmargin=2.2em]
\item $\srk_r(A)$ is nonincreasing in $r$, $\srk_r(A) = \srk_r(-A) \le \rank(A)$, and $\srk_r(A) = 0$ if $\max_{i,j}\|A_{ij}\| \le r$.
\item $\srk_r(A + B) \le \srk_r(A) + \srk_r(B)$.
\item $\srk_{rs}(AB) \le \srk_r(A) + \srk_s(B)$. In particular, if $U, W$ are bounded then $\srk_r(UAW) \le \srk_r(A)$, with equality if $U, W$ are unimodular.
\end{enumerate}
\end{lemma}

\begin{proof}
Part~(a) is immediate from the definition. For~(b), pick $A', B'$ attaining $\srk_r(A)$ and $\srk_r(B)$; by the ultrametric inequality, $\|(A+B)_{ij} - (A'+B')_{ij}\| \le \max\{\|A_{ij} - A'_{ij}\|, \|B_{ij} - B'_{ij}\|\} \le r$, so $\srk_r(A+B) \le \rank(A'+B') \le \rank(A') + \rank(B')$. For~(c), pick $A', B'$ attaining $\srk_r(A)$ and $\srk_s(B)$, and write
\[
AB - \bigl(A'B + (A - A')B'\bigr) = (A - A')(B - B').
\]
By the ultrametric inequality, every entry of $(A - A')(B - B')$ has norm at most $rs$, while $\rank(A'B + (A - A')B') \le \rank(A') + \rank(B')$. The last assertion follows by taking $s = 1$ (since $\srk_1(W) = 0$ for bounded $W$) and applying the inverse when $U, W$ are unimodular.
\end{proof}

\begin{proposition}[Smith normal form; Lean: \code{hasSmith\_all}, \code{srk\_diagonal}, \code{exists\_qinv}]
\label{prop:smith}
Let $A, B \in M_n(K)$ and $r \ge 0$.
\begin{enumerate}[label=\textup{(\alph*)},leftmargin=2.2em]
\item There exist unimodular matrices $U, W \in \GL_n(\OO_K)$ and a vector $d \in K^n$ such that $A = U \diag(d) W$.
\item $\srk_r(\diag(d)) = \#\{i : \|d_i\| > r\}$.
\item If $r > 0$, there exists $Q \in M_n(K)$ such that $\|(AQA - A)_{ij}\| \le r$ for all $i, j$, and for every $\rho \ge r$,
\[
\srk_{\rho^{-1}}(Q) + \srk_\rho(A) \le \srk_r(A).
\]
\item $\srk_r\bigl(\begin{smallmatrix} A & 0 \\ 0 & B \end{smallmatrix}\bigr) = \srk_r(A) + \srk_r(B)$.
\end{enumerate}
\end{proposition}

\begin{proof}
(a) Induct on $n$. If $A = 0$, take $U = W = I$ and $d = 0$. Otherwise, permute rows and columns (which are unimodular operations) so that $a := A_{11}$ has maximal norm among all entries of $A$. Writing $A = \bigl(\begin{smallmatrix} a & v^\mathsf{T} \\ u & C \end{smallmatrix}\bigr)$, the vectors $u a^{-1}$ and $a^{-1} v^\mathsf{T}$ are bounded, so the elementary triangular matrices $\bigl(\begin{smallmatrix} 1 & 0 \\ u a^{-1} & I \end{smallmatrix}\bigr)$ and $\bigl(\begin{smallmatrix} 1 & a^{-1} v^\mathsf{T} \\ 0 & I \end{smallmatrix}\bigr)$ are unimodular, and
\[
A = \begin{pmatrix} 1 & 0 \\ u a^{-1} & I \end{pmatrix}
\begin{pmatrix} a & 0 \\ 0 & C - u a^{-1} v^\mathsf{T} \end{pmatrix}
\begin{pmatrix} 1 & a^{-1} v^\mathsf{T} \\ 0 & I \end{pmatrix}.
\]
Applying the induction hypothesis to the $(n-1) \times (n-1)$ Schur complement $C - u a^{-1} v^\mathsf{T}$ finishes the proof.

(b) Let $I := \{i : \|d_i\| > r\}$. Setting $d'_i = d_i$ for $i \in I$ and $0$ otherwise gives $\max_{i,j}\|(\diag(d) - \diag(d'))_{ij}\| \le r$ and $\rank(\diag(d')) = |I|$, so $\srk_r(\diag(d)) \le |I|$. Conversely, let $A' \in M_n(K)$ satisfy $\|(\diag(d) - A')_{ij}\| \le r$ for all $i, j$, and let $M$ be the principal submatrix of $A'$ on $I \times I$. If $Mv = 0$ for some nonzero $v \in K^I$, choose $i_0 \in I$ maximizing $\|v_{i_0}\| > 0$. Then
\[
d_{i_0} v_{i_0} = \sum_{j \in I} (\diag(d) - A')_{i_0, j} v_j,
\]
so by the ultrametric inequality $\|d_{i_0}\|\|v_{i_0}\| \le r\|v_{i_0}\|$, contradicting $\|d_{i_0}\| > r$. Thus $M$ is invertible, and $\rank(A') \ge \rank(M) = |I|$.

(c) Write $A = U \diag(d) W$ as in~(a), define $q_i := d_i^{-1}$ if $\|d_i\| > r$ and $q_i := 0$ otherwise, and put $Q := W^{-1} \diag(q) U^{-1}$. Then $AQA - A = U \diag(d_i q_i d_i - d_i) W$, where each diagonal entry is $0$ (if $\|d_i\| > r$) or $-d_i$ (if $\|d_i\| \le r$), hence bounded in norm by $r$. Since $U, W$ are bounded, $\|(AQA - A)_{ij}\| \le r$ for all $i, j$. Moreover, for any $\rho \ge r$, by Lemma~\ref{lem:srk-basic}(c) and part~(b),
\[
\srk_{\rho^{-1}}(Q) = \#\{i : r < \|d_i\| < \rho\}, \qquad \srk_\rho(A) = \#\{i : \|d_i\| > \rho\}.
\]
These two subsets of $\{i : \|d_i\| > r\}$ are disjoint, so their cardinalities sum to at most $\#\{i : \|d_i\| > r\} = \srk_r(A)$.

(d) Apply~(a) to $A$ and $B$, form block-diagonal unimodular factors, and apply Lemma~\ref{lem:srk-basic}(c) and~(b).
\end{proof}

\subsection{Tight sequences modulo null sequences}

Let $(\iota_k)_{k \in \N}$ be a sequence of finite nonempty index sets, and let $\omega$ be a nonprincipal ultrafilter on $\N$. Write $\mathrm{Seq}(\iota, K) := \prod_{k \in \N} M_{\iota_k}(K)$, and for $A = (A_k) \in \mathrm{Seq}(\iota, K)$ and $r \ge 0$ define the normalized scale-$r$ rank
\[
u_r(A, k) := \frac{\srk_r(A_k)}{|\iota_k|} \in [0, 1].
\]
Call a sequence $A \in \mathrm{Seq}(\iota, K)$ \emph{tight} if
\[
\forall \varepsilon > 0,\ \exists \Lambda > 0,\ \forall^\omega k,\quad u_\Lambda(A, k) \le \varepsilon,
\]
and \emph{null} if
\[
\forall \varepsilon > 0,\ \forall r > 0,\ \forall^\omega k,\quad u_r(A, k) \le \varepsilon,
\]
where $\forall^\omega k$ means ``for $\omega$-almost every $k \in \N$'' (Lean: \code{IsTight}, \code{IsNull}). Let $\TT \subseteq \mathrm{Seq}(\iota, K)$ denote the set of tight sequences and $\NN \subseteq \TT$ the set of null sequences.

By Lemma~\ref{lem:srk-basic}, $\TT$ is a unital $K$-subalgebra of $\mathrm{Seq}(\iota, K)$ (note that every bounded sequence is tight with $\Lambda = 1$, and $c \cdot I$ is tight with $\Lambda = \max\{1, \|c\|\}$), and $\NN$ is a two-sided ideal of $\TT$: if $A \in \TT$ and $B \in \NN$, then given $\varepsilon, r > 0$, choosing $\Lambda > 0$ with $u_\Lambda(A, k) \le \varepsilon/2$ $\omega$-a.e.\ gives both $u_r(AB, k) \le u_\Lambda(A, k) + u_{r/\Lambda}(B, k) \le \varepsilon$ and $u_r(BA, k) \le u_{r/\Lambda}(B, k) + u_\Lambda(A, k) \le \varepsilon$ $\omega$-a.e.\ (Lean: \code{tight}, \code{nullCon}). Define the quotient $K$-algebra
\[
\RR(K, \iota, \omega) := \TT / \NN
\]
(Lean: \code{RankRings.RR}), and for $A \in \TT$ define (Lean: \code{rkS}, \code{rk})
\[
\rk([A]) := \inf\bigl\{c \ge 0 : \forall r > 0,\ \forall^\omega k,\ u_r(A, k) \le c\bigr\} = \sup_{r > 0} \lim_{k \to \omega} u_r(A, k) \in [0, 1].
\]

\begin{theorem}[Lean: \code{isRankFunction\_RR}, \code{rr\_regular}]
\label{thm:RR-regular}
The quotient $\RR(K, \iota, \omega)$ is a von Neumann regular $K$-algebra, and $\rk\colon \RR(K, \iota, \omega) \to [0, 1]$ is a faithful rank function.
\end{theorem}

\begin{proof}
First we check the rank function axioms (R1)--(R4). Clearly $0 \le \rk([A]) \le 1$, and $\rk([A]) = 0 \iff A \in \NN \iff [A] = 0$. Subadditivity $\rk([A+B]) \le \rk([A]) + \rk([B])$ and well-definedness modulo $\NN$ follow directly from Lemma~\ref{lem:srk-basic}(b). For $A, B \in \TT$ and $\varepsilon > 0$, picking $\Lambda > 0$ with $u_\Lambda(B, k) \le \varepsilon/2$ $\omega$-a.e.\ yields, for every $r > 0$,
\[
u_r(AB, k) \le u_{r/\Lambda}(A, k) + u_\Lambda(B, k) \le \rk([A]) + \varepsilon \qquad (\omega\text{-a.e.}),
\]
so $\rk([AB]) \le \rk([A])$; symmetrically $\rk([AB]) \le \rk([B])$. For $A = 1$, Proposition~\ref{prop:smith}(b) gives $u_{1/2}(1, k) = 1$ for all $k$, so $\rk(1) = 1$.

To verify~(R4) (additivity on orthogonal idempotents; Lean: \code{srk\_orth}, \code{rk\_add\_orth}), let $e = [E]$ and $f = [F]$ be orthogonal idempotents in $\RR(K, \iota, \omega)$ with $E, F \in \TT$. Consider the $2 \times 2$ block identity (at each coordinate $k$):
\begin{align*}
\begin{pmatrix} E_k & 0 \\ 0 & F_k \end{pmatrix}
&= \begin{pmatrix} E_k & 0 \\ F_k & 0 \end{pmatrix}
\begin{pmatrix} E_k + F_k & 0 \\ 0 & 0 \end{pmatrix}
\begin{pmatrix} E_k & F_k \\ 0 & 0 \end{pmatrix} \\
&\quad - \begin{pmatrix}
E_k(E_k+F_k)E_k - E_k & E_k(E_k+F_k)F_k \\
F_k(E_k+F_k)E_k & F_k(E_k+F_k)F_k - F_k
\end{pmatrix}.
\end{align*}
Each of the four blocks in the error matrix $D_k$ represents $0$ in $\RR(K, \iota, \omega)$ (since $e(e+f)e = e$, $e(e+f)f = 0$, $f(e+f)e = 0$, $f(e+f)f = f$), hence lies in $\NN$. For any $r, \Lambda > 0$, Proposition~\ref{prop:smith}(d) and Lemma~\ref{lem:srk-basic}(b),(c) (together with block subadditivity $\srk_r\bigl(\begin{smallmatrix} X & Y \\ Z & W \end{smallmatrix}\bigr) \le \srk_r(X) + \srk_r(Y) + \srk_r(Z) + \srk_r(W)$, which follows by writing a $2 \times 2$ block matrix as a sum of four bounded injections) yield
\[
u_r(E, k) + u_r(F, k) \le 2\bigl(u_\Lambda(E, k) + u_\Lambda(F, k)\bigr) + u_{r/\Lambda^2}(E + F, k) + u_r^{\mathrm{err}}(k),
\]
where $u_r^{\mathrm{err}}(k)$ is the sum of the $u_r$-values of the four error blocks. Given $\delta > 0$, choose $r > 0$ so that $u_r(E, k) > \rk(e) - \delta$ and $u_r(F, k) > \rk(f) - \delta$ $\omega$-a.e.; choose $\Lambda > 0$ so that $u_\Lambda(E, k) \le \delta$ and $u_\Lambda(F, k) \le \delta$ $\omega$-a.e. Since $u_{r/\Lambda^2}(E + F, k) \le \rk(e + f) + \delta$ and $u_r^{\mathrm{err}}(k) \le 4\delta$ $\omega$-a.e., taking a coordinate $k$ where all inequalities hold gives $\rk(e) + \rk(f) \le \rk(e + f) + 11\delta$. Letting $\delta \to 0$ gives $\rk(e + f) \ge \rk(e) + \rk(f)$, and hence equality by~(R2).

Finally, we prove von Neumann regularity (\code{rr\_regular}). Let $A \in \TT$ with $\rho := \rk([A])$. For each $j \in \N$, by definition of $\rho$, the property
\[
P(k, j) :\iff u_{1/(j+1)}(A, k) \le \rho + \frac{1}{j + 1}
\]
holds for $\omega$-a.e.\ $k$. Notice $P(k, 0)$ holds for every $k$ since $u_1(A, k) \le 1 \le \rho + 1$. Define
\[
J(k) := \max\{j \le k : P(k, j)\}, \qquad r_k := \frac{1}{J(k) + 1} > 0.
\]
Since $\omega$ is nonprincipal, for every $j \in \N$ both $j \le k$ and $P(k, j)$ hold $\omega$-a.e., so $J(k) \ge j$ $\omega$-a.e.; hence $\lim_{k \to \omega} r_k = 0$. For each $k$, let $B_k \in M_{\iota_k}(K)$ be the scale-$r_k$ quasi-inverse of $A_k$ provided by Proposition~\ref{prop:smith}(c).

To see that $B = (B_k)$ is tight, let $\varepsilon > 0$. By definition of $\rho$, there exists $r_0 > 0$ such that $u_{r_0}(A, k) > \rho - \varepsilon/2$ $\omega$-a.e. For $\omega$-a.e.\ $k$, we also have $r_k \le \min\{r_0, \varepsilon/2\}$. At every such $k$, Proposition~\ref{prop:smith}(c) and $P(k, J(k))$ give
\[
u_{r_0^{-1}}(B, k) \le u_{r_k}(A, k) - u_{r_0}(A, k) \le (\rho + r_k) - \Bigl(\rho - \frac{\varepsilon}{2}\Bigr) \le \varepsilon.
\]
Thus $B \in \TT$. Moreover, for every $r > 0$, we have $r_k \le r$ $\omega$-a.e., at which coordinates $\|(A_k B_k A_k - A_k)_{ij}\| \le r_k \le r$ for all $i, j$, so $u_r(ABA - A, k) = 0$. Hence $ABA - A \in \NN$, i.e., $[A][B][A] = [A]$ in $\RR(K, \iota, \omega)$.
\end{proof}

\subsection{The presentations \texorpdfstring{$B_\ell^\sigma$}{B\_l\^{}sigma} and \texorpdfstring{$\widetilde{B}_\ell^\sigma$}{Btilde\_l\^{}sigma}, adapted bases, and word combinatorics}

Let $\ell$ be a prime and $(\PP, \LL)$ be the projective plane $\mathbb{P}^2(\F_\ell)$, where $\PP = \LL = \mathbb{P}(\F_\ell^3)$ and a point $p = [v_p]$ lies on a line $l = [u_l]$ if and only if $u_l \cdot v_p = 0$.

\begin{lemma}[Lean: \code{RankRings.exists\_sigma}]
\label{lem:hall-sigma}
In any finite projective plane $(\PP, \LL)$ of order $q \ge 2$, there exists a bijection $\sigma\colon \PP \xrightarrow{\sim} \LL$ such that $p \notin \sigma(p)$ for all $p \in \PP$.
\end{lemma}

\begin{proof}
Consider the bipartite graph between $\PP$ and $\LL$ with an edge $(p, l)$ whenever $p \notin l$. Every vertex has degree $d := n - (q + 1) = q^2 > 0$. For any subset $S \subseteq \PP$ with neighborhood $N(S) \subseteq \LL$, double-counting edges between $S$ and $N(S)$ gives $d|S| \le d|N(S)|$, hence $|S| \le |N(S)|$. By Hall's marriage theorem and $|\PP| = |\LL|$, there is a bijection $\sigma\colon \PP \xrightarrow{\sim} \LL$ with $p \notin \sigma(p)$ for all $p \in \PP$.
\end{proof}

Fix such a $\sigma$ and $E = \{(p, p') : p \in \sigma(p')\}$. Alongside the presentation $B_\ell^\sigma$ in~\eqref{eq:B-ell}, consider the presentation $\widetilde{B}_\ell^\sigma$ (called a \emph{configuration}, Lean: \code{RankRings.Config}) with generators $f_p$ ($p \in \PP$) and $g_e, x_{e,i}, y_{e,i}$ ($e \in E$, $i \in \{0,1,2\}$) subject to
\begin{equation}
\label{eq:B-tilde}
f_p^2 = f_p, \qquad f_p g_e = g_e, \qquad f_{p'} g_e = 0, \qquad \sum_{i=0}^{2} y_{e,i} g_e x_{e,i} = 1 \qquad (e = (p, p') \in E).
\end{equation}
(Notice that the relations~\eqref{eq:B-ell} and~\eqref{eq:B-tilde} make sense for any finite projective plane $(\PP, \LL)$ and any bijection $\sigma\colon \PP \xrightarrow{\sim} \LL$.)

\begin{lemma}[Lean: \code{exists\_left\_id}, \code{config\_of\_relP}, \code{relP\_of\_config}]
\label{lem:B-equiv}
\leavevmode
\begin{enumerate}[label=\textup{(\alph*)},leftmargin=2.2em]
\item In any ring $R$, a solution of~\eqref{eq:B-tilde} yields a solution of~\eqref{eq:B-ell} upon replacing $x_{e,i}$ by $g_e x_{e,i}$.
\item Conversely, in any von Neumann regular ring $R$, a solution of~\eqref{eq:B-ell} yields a solution of~\eqref{eq:B-tilde}.
\end{enumerate}
\end{lemma}

\begin{proof}
Part~(a) is immediate. For~(b), it suffices to show by induction on $m \ge 0$ that for any elements $u, v, z_1, \dots, z_m \in R$ satisfying $u z_i = z_i$ and $v z_i = 0$ for all $i$, there exists $g \in R$ with $ug = g$, $vg = 0$, and $g z_i = z_i$ for all $i$. For $m = 0$, take $g = 0$. Given $g_0$ working for $z_1, \dots, z_{m-1}$, put $b := z_m - g_0 z_m$ and pick $q \in R$ with $bqb = b$. Then $ub = b$ and $vb = 0$. Setting $g := g_0 + bq - bq g_0$, we have $ug = g$, $vg = 0$, $g z_i = g_0 z_i + bq(z_i - g_0 z_i) = z_i$ for $i < m$, and $g z_m = g_0 z_m + bq b = g_0 z_m + b = z_m$. Applying this to $u = f_p$, $v = f_{p'}$, and $\{x_{e,0}, x_{e,1}, x_{e,2}\}$ for each $e = (p, p') \in E$ gives $g_e$ satisfying~\eqref{eq:B-tilde}.
\end{proof}

We now recall the word-counting lemmas needed for the tensor-power construction. For $m \ge 1$, let $W_m := \{0, 1, 2\}^m$ be the set of words of length $m$ over $\{0, 1, 2\} = \Z/3\Z$. For $I \subseteq \{0,1,2\}$ and $a \in \{0,1,2\}$, put $\cnt_I(w) := \#\{s : w_s \in I\}$ and $\cnt_a(w) := \cnt_{\{a\}}(w)$. Given a commutative ring $R$ and $C \in M_3(R)$, its $m$-th tensor power $C^{\otimes m} \in M_{W_m}(R)$ is indexed by words via $(C^{\otimes m})_{w, w'} = \prod_{s=1}^m C_{w_s, w'_s}$ (Lean: \code{tpow}). For $X \subseteq W_m$, write $\diag(X) \in M_{W_m}(R)$ for the diagonal $0/1$-matrix supported on $X$ (Lean: \code{dg}).

For $m = 3k + 1$, define the subsets of $W_m$:
\[
D_0(k) := \{w \in W_m : \cnt_0(w) \ge k + 1\}, \quad
T(m, k) := \{w \in W_m : \cnt_0(w) \ge k + 1,\ \cnt_2(w) \le k\}.
\]
For $i \in \Z/3\Z$, let $\operatorname{shift}_i(w) := (w_s - i \bmod 3)_{s=1}^m$, and define $M_i \in M_{W_m}(R)$ by $(M_i)_{v, w} = 1$ if $v = \operatorname{shift}_i(w) \in T(m, k)$ and $0$ otherwise (Lean: \code{T}, \code{shift}, \code{M}).

\begin{lemma}[Exact-third partition; Lean: \code{Words.lean}]
\label{lem:exact-third}
For $m = 3k + 1$ and every $w \in W_m$, there is a unique $i \in \Z/3\Z$ such that $\operatorname{shift}_i(w) \in T(m, k)$ (\code{existsUnique\_shift\_mem\_T}). Consequently, $\diag(T(m, k)) M_i = M_i$ for each $i$ and $\sum_{i \in \Z/3\Z} M_i^\mathsf{T} M_i = I$ (\code{sum\_transpose\_M\_mul\_M}).
\end{lemma}

\begin{proof}
Since $\cnt_a(\operatorname{shift}_i(w)) = \cnt_{a + i}(w)$, we have $\operatorname{shift}_i(w) \in T(m, k)$ if and only if $\cnt_i(w) \ge k + 1$ and $\cnt_{i - 1}(w) \le k$. Because $\cnt_0(w) + \cnt_1(w) + \cnt_2(w) = 3k + 1$, the three counts cannot all be $\le k$ (sum $\le 3k$) and cannot all be $\ge k + 1$ (sum $\ge 3k + 3$). Thus the cyclic sequence of truth values of $(\cnt_i(w) \ge k + 1)$ over $i \in \Z/3\Z$ is nonconstant and has length $3$, so it makes a unique transition from false to true: there is a unique $i \in \Z/3\Z$ with $\cnt_{i-1}(w) \le k$ and $\cnt_i(w) \ge k + 1$.

The identity $\diag(T(m, k)) M_i = M_i$ is immediate from the definition of $M_i$. Since $\operatorname{shift}_i$ is bijective, $(M_i^\mathsf{T} M_i)_{w, w'}$ equals $1$ if $w = w'$ and $\operatorname{shift}_i(w) \in T(m, k)$, and $0$ otherwise. Summing over $i \in \Z/3\Z$ gives $\delta_{w, w'}$ by uniqueness of $i$.
\end{proof}

\begin{lemma}[Off-diagonal blocks and anticoncentration; Lean: \code{offblock}, \code{window\_le}]
\label{lem:offblock-anti}
\leavevmode
\begin{enumerate}[label=\textup{(\alph*)},leftmargin=2.2em]
\item Let $R$ be a commutative ring, $\pi \in R$, $I, J \subseteq \{0,1,2\}$, and $C \in M_3(R)$ satisfy $\pi \mid C_{i,j}$ whenever $j \in I$ and $i \notin J$. Let $X, Y \subseteq W_m$ and $r \in \N$ satisfy $\cnt_J(w) < r$ for all $w \in X$ and $\cnt_I(w') \ge r$ for all $w' \in Y$. Then for every integer $t \ge 0$, there exist $Q_0, U_0 \in M_{W_m}(R)$ such that
\[
\diag(X)\, C^{\otimes m}\, \diag(Y) = \pi^{t+1} Q_0 + U_0 \diag\bigl(\{w' \in Y : \cnt_I(w') < r + t\}\bigr).
\]
\item For every integer $t \ge 0$ and every $\delta > 0$, there exists $K_0 \in \N$ such that for all $k \ge K_0$ and all $a \in \{0, 1, 2\}$,
\[
\#\bigl\{w \in W_{3k+1} : k - t \le \cnt_a(w) \le k + t\bigr\} \le \delta \cdot 3^{3k+1}.
\]
\end{enumerate}
\end{lemma}

\begin{proof}
(a) Split $Y = Y_{\mathrm{big}} \sqcup Y_{\mathrm{win}}$ where $Y_{\mathrm{win}} := \{w' \in Y : \cnt_I(w') < r + t\}$. Take $U_0 := \diag(X)\, C^{\otimes m}$. For $w \in X$ and $w' \in Y_{\mathrm{big}}$, let $A := \{s : w'_s \in I\}$ and $B := \{s : w_s \in J\}$; since $r \in \N$, we have $|A| \ge r + t$ and $|B| \le r - 1$, so $|A \setminus B| \ge |A| - |B| \ge t + 1$. For each $s \in A \setminus B$, we have $w'_s \in I$ and $w_s \notin J$, so $\pi \mid C_{w_s, w'_s}$. Hence $\pi^{t+1} \mid \prod_{s=1}^m C_{w_s, w'_s} = (C^{\otimes m})_{w, w'}$, yielding $\diag(X) C^{\otimes m} \diag(Y_{\mathrm{big}}) = \pi^{t+1} Q_0$ for some $Q_0 \in M_{W_m}(R)$.

(b) Put $m = 3k + 1$ and $N_m(c) := \#\{w \in W_m : \cnt_a(w) = c\} = 2^{m-c}\binom{m}{c}$. For any $B \ge 0$ with $3B + 1 \le 2k$ and any $c \le k + B$, we have $m - c \ge 2k + 1 - B$ and $2(c + 1) \le 2(k + B + 1)$, whence
\[
(2k - (3B + 1))\cdot 2(c + 1)
\le (2k - 3B - 1)(2k + 2B + 2)
= 2k(2k + 1 - B) - (6B^2 + 8B + 2)
\le 2k(m - c).
\]
Since $N_m(c+1) \cdot 2(c+1) = N_m(c) \cdot (m - c)$, this gives
\[
\bigl(2k - (3B + 1)\bigr) N_m(c) \le 2k\, N_m(c + 1).
\]
Iterating $r$ steps (by induction on $r$, or Bernoulli's inequality) whenever $c + r \le k + B + 1$ and $r(3B + 1) \le k$ yields
\[
\bigl(2k - r(3B + 1)\bigr) N_m(c) \le 2k\, N_m(c + r),
\qquad\text{hence}\qquad
N_m(c) \le 2 N_m(c + r).
\]
Now choose an integer $J > 2(2t + 1)/\delta$ and set $K_0 := J(3(t + J) + 1) + 1$. For every $k \ge K_0$ and every $c \in [k - t, k + t]$, taking $B = t + J$ gives $N_m(c) \le 2 N_m(c + r)$ for all $0 \le r \le J$. Summing over $r = 0, \dots, J$ yields
\[
(J + 1) N_m(c) \le 2\sum_{r=0}^J N_m(c + r) \le 2 \cdot 3^m.
\]
Summing over the at most $2t + 1$ integers $c \in [k - t, k + t]$ gives at most $\frac{2(2t+1)}{J+1} \cdot 3^m \le \delta \cdot 3^m$.
\end{proof}

\subsection{The exotic configuration in \texorpdfstring{$\RR(K, \iota, \omega)$}{RR(K, iota, omega)} and exact solutions over \texorpdfstring{$\F_\ell$}{F\_l}}
\label{sec:exotic}

For a prime $\ell$, write $\iota_k := W_{3k+1}$ (of cardinality $3^{3k+1}$; Lean: \code{RankRings.WI}). Let $\Z_\ell \subset \Q_\ell$ be the ring of $\ell$-adic integers, with reduction homomorphism $\Z_\ell \to \F_\ell = \Z/\ell\Z$.

\begin{theorem}[Lean: \code{RankRings.exotic}, \code{RankRings.exotic\_RR}]
\label{thm:exotic}
Let $\ell$ be a prime, let $K$ be an ultrametric normed field containing $\Q_\ell$ isometrically, let $\omega$ be a nonprincipal ultrafilter on $\N$, and let $\sigma\colon \PP \xrightarrow{\sim} \LL$ be a bijection of $\mathbb{P}^2(\F_\ell)$ with $p \notin \sigma(p)$ for all $p \in \PP$. Then $\RR(K, \iota, \omega)$ contains a solution of the configuration relations~\eqref{eq:B-tilde} (and hence of $B_\ell^\sigma$ in~\eqref{eq:B-ell}).
\end{theorem}

\begin{proof}
For each point $p \in \PP$, choose a nonzero representative $v_p \in \F_\ell^3$; for each line $l \in \LL$, choose a nonzero representative $u_l \in \F_\ell^3$ so that $p \in l \iff u_l \cdot v_p = 0$.

\emph{Step 1: Adapted $\ell$-adic bases} (\code{Geom.lean}, \code{Exotic.lean}: \code{exists\_B}, \code{exists\_A}).
For each $p \in \PP$, since $p \notin \sigma(p)$ we have $u_{\sigma(p)} \cdot v_p \neq 0$. Extending $v_p$ to an invertible matrix over $\F_\ell$ and subtracting multiples of column $0$ from columns $1$ and $2$, we obtain $\beta_p \in \GL_3(\F_\ell)$ whose column $0$ is $v_p$ and whose columns $1, 2$ are orthogonal to $u_{\sigma(p)}$ (\code{exists\_beta}). For each incident pair $(p, p') \in E$, since $u_{\sigma(p')} \cdot v_p = 0$ and $v_p, u_{\sigma(p')} \neq 0$, the orthogonal complement $u_{\sigma(p')}^\perp \subset \F_\ell^3$ is $2$-dimensional and contains $v_p$; choosing a vector $w \in u_{\sigma(p')}^\perp$ linearly independent from $v_p$ and extending $(v_p, w)$ to a basis of $\F_\ell^3$ yields $\alpha_{p, p'} \in \GL_3(\F_\ell)$ whose column $0$ is $v_p$ and whose column $1$ is orthogonal to $u_{\sigma(p')}$ (\code{exists\_alpha}). Lift each $\beta_p$ and $\alpha_{p,p'}$ entrywise to matrices $B_p, A_{p,p'} \in M_3(\Z_\ell)$; because $\det B_p \equiv \det \beta_p \not\equiv 0 \pmod \ell$ and $\det A_{p,p'} \equiv \det \alpha_{p,p'} \not\equiv 0 \pmod \ell$, both $B_p$ and $A_{p,p'}$ belong to $\GL_3(\Z_\ell)$ with reductions $\beta_p, \alpha_{p,p'}$ (\code{exists\_lift}).

We record two key congruence properties modulo $\ell$ for every $(p, p') \in E$ (\code{col0}, \code{row0}):
\begin{enumerate}[label=(\roman*),leftmargin=2.2em]
\item Since $\beta_p$ and $\alpha_{p,p'}$ share column $0$ (namely $v_p$), column $0$ of $\beta_p^{-1}\alpha_{p,p'}$ is $e_0$; hence
\[
\ell \mid (B_p^{-1} A_{p, p'})_{i, 0} \qquad\text{for all } i \in \{1, 2\}.
\]
\item For $j \in \{0, 1\}$, let $y \in \F_\ell^3$ be column $j$ of $\alpha_{p,p'}$; then $u_{\sigma(p')} \cdot y = 0$. Writing $y = \beta_{p'} (\beta_{p'}^{-1} y)$ and noting that $u_{\sigma(p')}^\mathsf{T} \beta_{p'} = (u_{\sigma(p')} \cdot v_{p'},\, 0,\, 0)$ with $u_{\sigma(p')} \cdot v_{p'} \neq 0$, we obtain $0 = u_{\sigma(p')} \cdot y = (u_{\sigma(p')} \cdot v_{p'})(\beta_{p'}^{-1} y)_0$, so $(\beta_{p'}^{-1} \alpha_{p,p'})_{0, j} = 0$. Hence
\[
\ell \mid (B_{p'}^{-1} A_{p, p'})_{0, j} \qquad\text{for all } j \in \{0, 1\}.
\]
\end{enumerate}

\emph{Step 2: Tensor-power sequences.}
For each $k \in \N$, put $m = 3k + 1$ and define matrices over $\Z_\ell$ indexed by $\iota_k = W_m$:
\begin{align*}
F_p^{(k)} &:= B_p^{\otimes m} \diag(D_0(k)) (B_p^{-1})^{\otimes m}, &
G_{p,p'}^{(k)} &:= A_{p,p'}^{\otimes m} \diag(T(m,k)) (A_{p,p'}^{-1})^{\otimes m}, \\
X_{p,p',i}^{(k)} &:= A_{p,p'}^{\otimes m} M_i, &
Y_{p,p',i}^{(k)} &:= M_i^\mathsf{T} (A_{p,p'}^{-1})^{\otimes m} \qquad (i \in \Z/3\Z).
\end{align*}
All entries lie in $\Z_\ell$, hence have norm at most $1$ in $K$, so each sequence defines a tight sequence in $\TT$ and thus an element $f_p, g_{p,p'}, x_{p,p',i}, y_{p,p',i} \in \RR(K, \iota, \omega)$.

For every $k \ge 0$, we have the exact matrix identities $(F_p^{(k)})^2 = F_p^{(k)}$ and, by Lemma~\ref{lem:exact-third},
\[
\sum_{i \in \Z/3\Z} Y_{p,p',i}^{(k)} G_{p,p'}^{(k)} X_{p,p',i}^{(k)}
= \sum_{i \in \Z/3\Z} M_i^\mathsf{T} \diag(T(m,k)) M_i
= \sum_{i \in \Z/3\Z} M_i^\mathsf{T} M_i = I.
\]
Hence $f_p^2 = f_p$ and $\sum_{i} y_{p,p',i} g_{p,p'} x_{p,p',i} = 1$ in $\RR(K, \iota, \omega)$.

\emph{Step 3: Nullity of the incidence errors.}
Fix $(p, p') \in E$, an integer $t \ge 0$, and $\delta > 0$. By Lemma~\ref{lem:offblock-anti}(b), choose $K_0$ so that the window of counts in $[k - t, k + t]$ has cardinality at most $\delta \cdot 3^{3k+1}$ for all $k \ge K_0$ and all letters $a \in \{0,1,2\}$. Fix $k \ge K_0$ and put $m = 3k + 1$.

First, for $(I - F_p^{(k)}) G_{p,p'}^{(k)}$: using $I - \diag(D_0(k)) = \diag(D_0(k)^c)$,
\[
(I - F_p^{(k)}) G_{p,p'}^{(k)}
= B_p^{\otimes m} \bigl(\diag(D_0(k)^c)\, (B_p^{-1} A_{p,p'})^{\otimes m}\, \diag(T(m,k))\bigr) (A_{p,p'}^{-1})^{\otimes m}.
\]
By Step~1(i), the matrix $C = B_p^{-1} A_{p,p'}$ satisfies $\ell \mid C_{i,j}$ for $j \in \{0\}$ and $i \notin \{0\}$. Every $w \in D_0(k)^c$ has $\cnt_{\{0\}}(w) < k + 1$, and every $w' \in T(m, k)$ has $\cnt_{\{0\}}(w') \ge k + 1$. Applying Lemma~\ref{lem:offblock-anti}(a) with $I = J = \{0\}$ and $r = k + 1$ yields $Q_0, U_0 \in M_{W_m}(\Z_\ell)$ and the window subset
\[
W := \{w' \in T(m, k) : \cnt_0(w') < k + 1 + t\} \subseteq \{w' \in W_m : k - t \le \cnt_0(w') \le k + t\}
\]
such that setting $Q := B_p^{\otimes m} Q_0 (A_{p,p'}^{-1})^{\otimes m}$, $U := B_p^{\otimes m} U_0$, and $Z := (A_{p,p'}^{-1})^{\otimes m}$ in $M_{W_m}(\Z_\ell)$ gives $(I - F_p^{(k)}) G_{p,p'}^{(k)} = \ell^{t+1} Q + U \diag(W) Z$ with $|W| \le \delta \cdot 3^{3k+1}$.

Second, for $F_{p'}^{(k)} G_{p,p'}^{(k)}$:
\[
F_{p'}^{(k)} G_{p,p'}^{(k)}
= B_{p'}^{\otimes m} \bigl(\diag(D_0(k))\, (B_{p'}^{-1} A_{p,p'})^{\otimes m}\, \diag(T(m,k))\bigr) (A_{p,p'}^{-1})^{\otimes m}.
\]
By Step~1(ii), $C = B_{p'}^{-1} A_{p,p'}$ satisfies $\ell \mid C_{i,j}$ for $j \in \{0, 1\}$ and $i \notin \{1, 2\}$ (i.e., $i = 0$). For every $w \in D_0(k)$, $\cnt_0(w) \ge k + 1$ implies $\cnt_{\{1,2\}}(w) = m - \cnt_0(w) < 2k + 1$; for every $w' \in T(m, k)$, $\cnt_2(w') \le k$ implies $\cnt_{\{0,1\}}(w') = m - \cnt_2(w') \ge 2k + 1$. Applying Lemma~\ref{lem:offblock-anti}(a) with $I = \{0,1\}$, $J = \{1,2\}$, and $r = 2k + 1$ yields $Q'_0, U'_0 \in M_{W_m}(\Z_\ell)$ and
\[
W' := \{w' \in T(m, k) : \cnt_{\{0,1\}}(w') < 2k + 1 + t\} \subseteq \{w' \in W_m : k - t \le \cnt_2(w') \le k + t\}
\]
such that setting $Q' := B_{p'}^{\otimes m} Q'_0 (A_{p,p'}^{-1})^{\otimes m}$, $U' := B_{p'}^{\otimes m} U'_0$, and $Z' := (A_{p,p'}^{-1})^{\otimes m}$ in $M_{W_m}(\Z_\ell)$ gives $F_{p'}^{(k)} G_{p,p'}^{(k)} = \ell^{t+1} Q' + U' \diag(W') Z'$ with $|W'| \le \delta \cdot 3^{3k+1}$.

Finally, to see that both error sequences lie in $\NN$ (\code{ExoticS.lean}: \code{isNull\_of\_approx}), let $D = (D^{(k)})$ be either sequence and let $\varepsilon, r > 0$. Since $\|\ell\| = \ell^{-1} < 1$ in $K$, choose $t \in \N$ with $\ell^{-(t+1)} \le r$. For all $k \ge K_0(t, \varepsilon)$, the decomposition $D^{(k)} = \ell^{t+1} Q + U \diag(W) Z$ has $\max_{i,j}\|(\ell^{t+1} Q)_{ij}\| \le \ell^{-(t+1)} \le r$ and $\rank(U \diag(W) Z) \le |W| \le \varepsilon \cdot 3^{3k+1}$, so $u_r(D, k) \le \varepsilon$. Since $\omega$ is nonprincipal, $k \ge K_0$ holds $\omega$-a.e., hence $(1 - f_p) g_{p,p'} = 0$ and $f_{p'} g_{p,p'} = 0$ in $\RR(K, \iota, \omega)$.
\end{proof}

\begin{proposition}[Exact representations over $\F_\ell$; Lean: \code{exists\_config\_modL}, \code{rel\_modL}]
\label{prop:modL}
For every prime $\ell$, every fixed-incidence-free bijection $\sigma$ of $\mathbb{P}^2(\F_\ell)$, and every integer $k \ge 0$, the relations~\eqref{eq:B-tilde} (and hence~\eqref{eq:B-ell} of $B_\ell^\sigma$) admit an exact solution in $M_{3^{3k+1}}(\F_\ell)$.
\end{proposition}

\begin{proof}
Perform Step~2 of the proof of Theorem~\ref{thm:exotic} directly over $\F_\ell$ using the matrices $\beta_p, \alpha_{p,p'} \in \GL_3(\F_\ell)$ in place of $B_p, A_{p,p'}$. In Step~1(i) and~1(ii), divisibility by $\ell$ becomes exact vanishing ($(\beta_p^{-1}\alpha_{p,p'})_{i,0} = 0$ for $i \neq 0$, and $(\beta_{p'}^{-1}\alpha_{p,p'})_{0,j} = 0$ for $j \in \{0,1\}$). Applying Lemma~\ref{lem:offblock-anti}(a) with $\pi = 0$ and $t = 0$ (so the window $\{w' \in Y : \cnt_I(w') < r\}$ is empty) forces both off-diagonal blocks to vanish identically (\code{ModL.lean}: \code{offblock\_zero}). Thus $(I - F_p^{(k)}) G_{p,p'}^{(k)} = 0$ and $F_{p'}^{(k)} G_{p,p'}^{(k)} = 0$ hold exactly in $M_{W_{3k+1}}(\F_\ell) \cong M_{3^{3k+1}}(\F_\ell)$.
\end{proof}

\section{Proofs of the matrix-ultraproduct obstruction theorems}
\label{sec:proofs}

We now assemble the results of Sections~\ref{sec:ineq} and~\ref{sec:ultrametric} to prove Theorems~\ref{thm:main-ultra} and~\ref{thm:finite-form}, corresponding to the Lean modules \code{NoHom.lean}, \code{NoHomF.lean}, \code{CAlg.lean}, \code{Main.lean}, \code{Final.lean}, \code{FiniteForm.lean}, \code{FiniteFormF.lean}, and \code{ModL.lean}.

\subsection{Exact lifting and no configurations in characteristic-zero rank ultraproducts}

Let $I$ be an index set, $\omega$ an ultrafilter on $I$, $(d_i)_{i \in I}$ positive integers, and $(F_i)_{i \in I}$ fields. Write $\mathcal{M} := \prod_{i \in I} M_{d_i}(F_i)$ and $J_\omega \subset \mathcal{M}$ for the rank-null ideal. Because every matrix over a field admits a quasi-inverse, there is a map $q\colon \mathcal{M} \to \mathcal{M}$ satisfying $A\, q(A)\, A = A$ for all $A \in \mathcal{M}$ (in particular, $\RUF{\omega}{d_i}{F_i} = \mathcal{M}/J_\omega$ is von Neumann regular).

\begin{lemma}[Exact lifting; Lean: \code{Gen.idem\_lift}, \code{Gen.edge\_lift}]
\label{lem:lifting}
\leavevmode
\begin{enumerate}[label=\textup{(\alph*)},leftmargin=2.2em]
\item If $A \in \mathcal{M}$ satisfies $A^2 - A \in J_\omega$, there exists $A' \in \mathcal{M}$ with $(A')^2 = A'$ and $A' - A \in J_\omega$.
\item If $F_1, F_2, G \in \mathcal{M}$ satisfy $F_1 G - G \in J_\omega$ and $F_2 G \in J_\omega$, there exists $a \in \mathcal{M}$ with $F_1 a = a$, $F_2 a = 0$, and $a - G \in J_\omega$.
\end{enumerate}
\end{lemma}

\begin{proof}
(a) Put $D := A^2 - A \in J_\omega$ and $e := 1 - q(D)D$. Then $De = D - D q(D) D = 0$ and $DA = AD$, so $eAe = Ae - q(D)DAe = Ae - q(D)ADe = Ae$. Setting $A' := Ae$, we obtain
\[
(A')^2 = A(eAe) = A^2 e = (D + A)e = Ae = A',
\]
and $A' - A = -A q(D) D \in J_\omega$.

(b) Put $h := F_1 G - G \in J_\omega$ and $\pi_1 := 1 - q(h)h$, so $h\pi_1 = 0$ and hence $F_1(G\pi_1) = (h + G)\pi_1 = G\pi_1$. Next put $h' := F_2 G \pi_1 \in J_\omega$ and $\pi_2 := 1 - q(h')h'$, so $h'\pi_2 = 0$. Setting $a := G\pi_1\pi_2$, we have $F_1 a = F_1(G\pi_1)\pi_2 = G\pi_1\pi_2 = a$, $F_2 a = h'\pi_2 = 0$, and $a - G = -G q(h)h\pi_2 - G q(h')h' \in J_\omega$.
\end{proof}

\begin{proposition}[Lean: \code{Gen.exists\_embed}, \code{Gen.no\_config}]
\label{prop:no-config}
Let $(\PP, \LL)$ be any finite projective plane of order $q \ge 7$ with a bijection $\sigma\colon \PP \xrightarrow{\sim} \LL$. For any index set $I$, ultrafilter $\omega$ on $I$, positive integers $(d_i)_{i \in I}$ ($d_i \ge 1$), and fields $(F_i)_{i \in I}$ of characteristic zero, the rank ultraproduct $\RUF{\omega}{d_i}{F_i}$ contains no solution of the configuration relations~\eqref{eq:B-tilde}.
\end{proposition}

\begin{proof}
Suppose for contradiction that $f_p, g_{p,p'}, x_{p,p',j}, y_{p,p',j} \in \mathcal{M}/J_\omega$ satisfy~\eqref{eq:B-tilde}. Lift them to sequences $F_p, G_{p,p'}, X_{p,p',j}, Y_{p,p',j} \in \mathcal{M}$. By Lemma~\ref{lem:lifting}(a), replace each $F_p$ by an exact idempotent $F'_p \in \mathcal{M}$ with $F'_p - F_p \in J_\omega$. For each $(p, p') \in E$, we still have $F'_p G_{p,p'} - G_{p,p'} \in J_\omega$ and $F'_{p'} G_{p,p'} \in J_\omega$, so Lemma~\ref{lem:lifting}(b) yields $a_{p,p'} \in \mathcal{M}$ with $F'_p a_{p,p'} = a_{p,p'}$, $F'_{p'} a_{p,p'} = 0$, and $a_{p,p'} - G_{p,p'} \in J_\omega$.

Set $r_{p,p'} := a_{p,p'}\, q(a_{p,p'}) \in \mathcal{M}$. Then $r_{p,p'}^2 = r_{p,p'}$, $r_{p,p'} a_{p,p'} = a_{p,p'}$, $F'_p r_{p,p'} = r_{p,p'}$, and $F'_{p'} r_{p,p'} = 0$. Define the error sequence
\[
\mathcal{E}_{p,p'} := \sum_{j=0}^{2} Y_{p,p',j} a_{p,p'} X_{p,p',j} - 1 \in J_\omega.
\]
Since $q \ge 7$, we have $q^2 - 7q + 1 > 0$, i.e., $9q < (q+1)^2$, so $\sqrt{q} < (q+1)/3$. Put
\[
\varepsilon := \frac{(q + 1)/3 - \sqrt{q}}{4(q + 1)} > 0.
\]
Because $E$ is finite and each $\mathcal{E}_{p,p'} \in J_\omega$, there exists a coordinate $i_0 \in I$ such that $\rank(\mathcal{E}_{p,p'}(i_0)) < \varepsilon d_{i_0}$ for all $(p, p') \in E$. At $i_0$, rewriting the identity
\[
I = \sum_{j=0}^2 \bigl(Y_{p,p',j}(i_0) r_{p,p'}(i_0)\bigr)\bigl(a_{p,p'}(i_0) X_{p,p',j}(i_0)\bigr) - \mathcal{E}_{p,p'}(i_0)
\]
and applying subadditivity and multiplicativity of matrix rank gives
\[
d_{i_0} \le 3\rank\bigl(r_{p,p'}(i_0)\bigr) + \varepsilon d_{i_0},
\qquad\text{so}\qquad
\rank\bigl(r_{p,p'}(i_0)\bigr) \ge \frac{1 - \varepsilon}{3}\, d_{i_0}
\]
for every $(p, p') \in E$.

Now let $S_0 \subset F_{i_0}$ be the finite set of all entries of the matrices $F'_p(i_0)$ ($p \in \PP$) and $r_{p,p'}(i_0)$ ($(p,p') \in E$), and let $K_0 := \Q(S_0) \subseteq F_{i_0}$ be the countable subfield they generate. Adjoining $\R$-many indeterminates to $K_0$ and passing to the algebraic closure of the resulting fraction field produces an algebraically closed field $L$ of characteristic zero and cardinality $\mathfrak{c} = |\C|$, which is isomorphic to $\C$ by Steinitz's theorem; thus there exists a field embedding $\varphi\colon K_0 \hookrightarrow \C$ (\code{Gen.exists\_embed}). Applying $\varphi$ entrywise to $F'_p(i_0)$ and $r_{p,p'}(i_0)$ produces matrices $\widetilde{F}_p, \widetilde{r}_{p,p'} \in M_{d_{i_0}}(\C)$ satisfying all hypotheses of Theorem~\ref{thm:ineq-C}. Moreover, because an idempotent matrix over any characteristic-zero field has trace equal to its rank times $1$, we have $\tr(\widetilde{r}_{p,p'}) = \rank(r_{p,p'}(i_0))$. Theorem~\ref{thm:ineq-C} therefore implies
\[
4(q + 1) n \cdot \frac{1 - \varepsilon}{3}\, d_{i_0}
\le 4\sum_{(p, p') \in E} \RePart\tr(\widetilde{r}_{p,p'})
\le (q + 1 + \sqrt{q})\, n d_{i_0}.
\]
Dividing by $n d_{i_0} > 0$ yields $\frac{4}{3}(q + 1)(1 - \varepsilon) \le q + 1 + \sqrt{q}$, i.e., $\frac{q+1}{3} - \sqrt{q} \le \frac{4}{3}(q+1)\varepsilon = \frac{1}{3}\bigl(\frac{q+1}{3} - \sqrt{q}\bigr)$, contradicting $\frac{q+1}{3} - \sqrt{q} > 0$.
\end{proof}

\subsection{Proofs of Theorems~\ref{thm:main-ultra} and~\ref{thm:finite-form}}

\begin{proof}[Proof of Theorem~\ref{thm:main-ultra}]
Fix $\ell = 7$, a fixed-incidence-free bijection $\sigma$ of $\mathbb{P}^2(\F_7)$ (Lemma~\ref{lem:hall-sigma}), and a nonprincipal ultrafilter $\omega$ on $\N$ (e.g., the hyperfilter \code{Filter.hyperfilter Nat}).

(a) Let $R_0 := \RR(\Q_7, \iota, \omega)$. By Theorems~\ref{thm:RR-regular} and~\ref{thm:exotic}, $R_0$ is a von Neumann regular $\Q_7$-algebra with faithful rank function $\rk$ containing a solution $c$ of~\eqref{eq:B-tilde}. Choose a function $\mathrm{qi}\colon R_0 \to R_0$ satisfying $a \cdot \mathrm{qi}(a) \cdot a = a$ for all $a \in R_0$, and let $S \subseteq R_0$ be the smallest subring containing $\Q \cdot 1_{R_0}$ and the finitely many elements of $c$ that is closed under $\mathrm{qi}$ (Lean: \code{RankRings.clos}, generated by the countable inductive term type \code{Tm}). Then $S$ is countable, von Neumann regular, a $\Q$-algebra (as $\Q \cdot 1_{R_0}$ is central in $R_0$), inherits the faithful rank function $\rk|_S$, and contains the configuration $c$ (\code{RankRings.main\_core}). If $\psi\colon S \to \RUF{\omega'}{d_i}{F_i}$ were any unital ring homomorphism into a characteristic-zero matrix rank ultraproduct, then $\psi(c)$ would be a solution of~\eqref{eq:B-tilde} in $\RUF{\omega'}{d_i}{F_i}$, contradicting Proposition~\ref{prop:no-config} (\code{main\_countable\_char0}, \code{elek\_countable}).

(b) Let $K := \overline{\Q}_7$ be the algebraic closure of $\Q_7$, equipped with the unique extension of the $7$-adic absolute value (Mathlib: \code{PadicAlgCl 7}). Since $\Q_7$ has cardinality $\mathfrak{c}$ and $K/\Q_7$ is algebraic, $K$ is an algebraically closed field of characteristic zero and cardinality $\mathfrak{c}$, so Steinitz's theorem (using the Axiom of Choice) provides a non-constructive ring isomorphism $e\colon K \xrightarrow{\sim} \C$ (\code{CAlg.lean}: \code{nonempty\_padicAlgCl\_equiv\_complex}). Setting $R := \RR(K, \iota, \omega)$ and pulling back the central $K$-algebra structure along $e^{-1}\colon \C \xrightarrow{\sim} K$ makes $R$ a von Neumann regular $\C$-algebra with faithful rank function $\rk$ containing a solution of~\eqref{eq:B-tilde}. By Proposition~\ref{prop:no-config}, $R$ admits no unital ring homomorphism into any characteristic-zero matrix rank ultraproduct (\code{main\_complex}, \code{elek\_complex}).
\end{proof}

\begin{proof}[Proof of Theorem~\ref{thm:finite-form}]
Part~(a) follows from Theorem~\ref{thm:main-ultra}(a) and Lemma~\ref{lem:B-equiv}(a), and part~(c) is Proposition~\ref{prop:modL} combined with Lemma~\ref{lem:B-equiv}(a).

For part~(b) (Lean: \code{eps\_char0\_paper}, \code{finite\_form\_paper}), suppose for contradiction that no such $\varepsilon > 0$ exists. Then for every $k \in \N$ there exist a field $F_k$ of characteristic zero, an integer $N_k \ge 1$, and matrices $A^{(k)}\colon \mathrm{Gn}_\sigma \to M_{N_k}(F_k)$ such that
\[
\rank\bigl(P_j(A^{(k)})\bigr) \le \frac{1}{k + 1}\, N_k \qquad\text{for every relation } P_j \text{ of } B_7^\sigma.
\]
Let $\omega$ be a nonprincipal ultrafilter on $\N$, and consider the images $a_g := [(A^{(k)}_g)_{k \in \N}]$ in the rank ultraproduct $\mathcal{Q} := \RUF{\omega}{N_k}{F_k}$. For every relation $P_j$, we have $\lim_{k \to \omega} \rank(P_j(A^{(k)}))/N_k \le \lim_{k \to \omega} (k+1)^{-1} = 0$, so $P_j(a) = 0$ in $\mathcal{Q}$. Because $\mathcal{Q}$ is von Neumann regular (\code{FiniteFormF.lean}: \code{regular\_RUF}), Lemma~\ref{lem:B-equiv}(b) upgrades $a$ to a solution of~\eqref{eq:B-tilde} in $\mathcal{Q}$, contradicting Proposition~\ref{prop:no-config}. (The exact same argument gives the uniform $\varepsilon > 0$ for the presentation $\widetilde{B}_7^\sigma$ directly, without invoking Lemma~\ref{lem:B-equiv}; Lean: \code{eps\_char0}, \code{finite\_form\_char0}.)
\end{proof}

\section{Operator-algebraic corollaries: finite von Neumann algebras and tracial \texorpdfstring{$C^*$}{C*}-algebras}
\label{sec:vn}

\noindent\textbf{Status note.} \emph{The results in this section are proved on paper and have not yet been formalized in Lean~4.}

\subsection{Preliminaries on affiliated operators}

Let $M \subseteq B(\HH)$ be a finite von Neumann algebra acting on a Hilbert space $\HH$, and let $\U(M)$ be the $*$-regular ring of closed, densely defined operators affiliated with $M$~\cite{MurrayVonNeumann,Berberian1957,Luck2002}. Suppose first that $M$ carries a faithful normal tracial state $\tau$ ($\tau(1) = 1$). We write $\|x\|_2 := \tau(x^*x)^{1/2}$ and $\langle x, y \rangle_\tau := \tau(y^* x)$ for $x, y \in M$. We use four standard properties of $\U(M)$ (see~\cite[Chapter~8]{Luck2002}):
\begin{enumerate}[label=(U\arabic*),leftmargin=2.8em]
\item Every $x \in \U(M)$ has a polar decomposition $x = v|x|$ with $v \in M$ a partial isometry, $s_r(x) := v^*v \in M$ (right support, onto $(\ker x)^\perp$), and $s_l(x) := vv^* \in M$ (left support, onto $\overline{\ran(x)}$). Setting $\phi(t) = t^{-1}$ for $t > 0$ and $\phi(0) = 0$, the element $c := \phi(|x|) v^* \in \U(M)$ satisfies $xc = s_l(x)$ and $cx = s_r(x)$ (as $|x|\phi(|x|) = \phi(|x|)|x| = s_r(x)$). In particular, $x\U(M) = s_l(x)\U(M)$; if $wx = x$ for a projection $w \in M$ then $s_l(x) \le w$, and if $xw = x$ then $s_r(x) \le w$.
\item Defining $\rk_\tau(x) := \tau(s_l(x)) = \tau(s_r(x))$, $\rk_\tau$ is a faithful rank function on $\U(M)$; in particular $|\rk_\tau(x) - \rk_\tau(y)| \le \rk_\tau(x - y)$. Each principal right ideal $e\U(M)$ ($e \in M$ a projection) determines $e$ uniquely; we set $\dim_\tau(e\U(M)) := \tau(e)$, and we have $e\U(M) \cap e'\U(M) = (e \wedge e')\U(M)$. For an idempotent $f \in \U(M)$, we define its range projection $r_f := s_l(f)$ and its kernel projection $k_f := s_l(1 - f)$, which satisfy $f\U(M) = r_f\U(M)$ and $(1 - f)\U(M) = k_f\U(M)$.
\item For any $a \in \U(M)$, the spectral projections of $|a|$ lie in $M$, the truncation $a\1_{[0,k]}(|a|) = v(|a|\1_{[0,k]}(|a|))$ lies in $M$ with norm at most $k$, and $\tau(\1_{(k,\infty)}(|a|)) \to 0$ as $k \to \infty$ by normality of $\tau$.
\item For invertible $g \in M$ and a projection $e \in M$, let $e^g := s_l(ge) \in M$ be the orthogonal projection onto the closed subspace $g(e\HH)$. Since $g$ is invertible, $\ker(ge) = \ker(e)$, so $s_r(ge) = e$ and $\tau(e^g) = \tau(s_l(ge)) = \tau(s_r(ge)) = \tau(e)$. Moreover, $(e \wedge e')^g = e^g \wedge (e')^g$ because $g$ is a bijection of $\HH$.
\end{enumerate}

\subsection{The incidence inequality in \texorpdfstring{$\U(M)$}{U(M)}}

\begin{theorem}[Paper proof, not formalized in Lean~4]
\label{thm:vn-ineq}
Let $(M, \tau)$ be a finite von Neumann algebra with a faithful normal tracial state, let $(\PP, \LL)$ be a finite projective plane of order $q \ge 2$ with $n = q^2 + q + 1$ points, let $\sigma\colon \PP \xrightarrow{\sim} \LL$ be a bijection, and set $E := \{(p, p') : p \in \sigma(p')\}$. Then for all idempotents $f_p \in \U(M)$ ($p \in \PP$), bounded or unbounded,
\begin{equation}
\label{eq:vn-ineq}
\sum_{(p, p') \in E} \dim_\tau\bigl(f_p \U(M) \cap (1 - f_{p'})\U(M)\bigr) \le \frac{n(q + 1 + \sqrt{q})}{4}.
\end{equation}
\end{theorem}

\begin{proof}
We first prove~\eqref{eq:vn-ineq} when every $f_p$ lies in $M$, and then deduce the general unbounded case by spectral truncation.

\emph{Step 1: Bounded idempotents and their canonical projections.}
For an idempotent $f \in M$ and $f' := 1 - f$, put $L(f) := f^*f + (f')^*f' \in M$. For every $\xi \in \HH$, $\|\xi\|^2 = \|f\xi + f'\xi\|^2 \le 2(\|f\xi\|^2 + \|f'\xi\|^2) = 2\langle L(f)\xi, \xi\rangle$, so $L(f) \ge \frac{1}{2}\cdot 1$ is positive and invertible in $M$. Put $P := f L(f)^{-1}$ and $Q := f' L(f)^{-1}$. Exactly as in Lemma~\ref{lem:L-idem}(a)--(c), $L(f)f = f^*f = f^* L(f)$ implies $P^* = P = P^2$, $Pf = f$, and $fP = P$, so $P = r_f$ and $Q = r_{f'} = k_f$ are orthogonal projections in $M$ satisfying
\begin{equation}
\label{eq:vn-PQ}
P + Q = L(f)^{-1}, \qquad f L(f)^{-1} f^* = P, \qquad f' L(f)^{-1} (f')^* = Q, \qquad \tau(P) + \tau(Q) = 1.
\end{equation}
(Here $f L(f)^{-1} f^* = P f^* = (fP)^* = P^* = P$, and $\tau(P) = \tau(fP) = \tau(Pf) = \tau(f)$.) Moreover, since $\log t \le t - 1$ for $t > 0$, continuous functional calculus gives $\log(L(f)^{-1}) \le L(f)^{-1} - 1$, hence $\tau(\log L(f)) = -\tau(\log L(f)^{-1}) \ge 1 - \tau(L(f)^{-1}) = 0$.

\emph{Step 2: Potential and first derivative.}
For invertible $g \in M$, put $E_p(g) := g f_p g^{-1} \in M$, let $P_p(g) := r_{E_p(g)} = (r_{f_p})^g$ and $Q_p(g) := k_{E_p(g)} = (k_{f_p})^g$, and define
\[
D(g) := \sum_{p \in \PP} \bigl(P_p(g) + Q_p(g) - 1\bigr) \in M_{\mathrm{sa}}, \qquad
\Phi(g) := \sum_{p \in \PP} \tau\bigl(\log L(E_p(g))\bigr) \ge 0.
\]
Fix $H = H^* \in M$ and $g \in \GL(M)$, and consider the smooth curve $g_t := e^{tH} g$ ($t \in \R$). Put $E_p(t) := E_p(g_t) = e^{tH} E_p(0) e^{-tH}$, $F_p(t) := 1 - E_p(t)$, and $L_p(t) := L(E_p(t))$. Differentiating in operator norm yields $E'_p = H E_p - E_p H$, $F'_p = H F_p - F_p H$, and
\[
L'_p = 2(E_p^* H E_p + F_p^* H F_p) - H L_p - L_p H.
\]
We claim that $\frac{d}{dt}\tau(\log L_p(t)) = \tau(L_p(t)^{-1} L'_p(t))$. For any positive $L \in M$ with $L \ge \frac{1}{2}$, continuous functional calculus on $\spn(L) \subseteq [\frac{1}{2}, \|L\|]$ gives the norm-convergent improper Riemann integral
\[
\log L = \int_0^\infty \bigl((1 + s)^{-1} - (L + s)^{-1}\bigr)\, ds.
\]
Put $R_t(s) := (L_p(t) + s)^{-1}$. Since $L_p(t) \ge \frac{1}{2}$, we have $\|R_t(s)\| \le (\frac{1}{2} + s)^{-1}$ for all $t \in \R$ and $s \ge 0$. For $0 < |h| \le 1$, the resolvent identity yields
\[
\frac{\tau(\log L_p(t+h)) - \tau(\log L_p(t))}{h}
= \int_0^\infty \tau\Bigl(R_{t+h}(s)\,\frac{L_p(t+h) - L_p(t)}{h}\, R_t(s)\Bigr)\, ds.
\]
As $h \to 0$, the integrand converges pointwise to $\tau(R_t(s) L'_p(t) R_t(s))$ and is uniformly dominated by $K / (\frac{1}{2} + s)^2 \in L^1([0, \infty))$, where $K := \sup_{|\theta| \le 1} \|L'_p(t + \theta)\| < \infty$. By Lebesgue's dominated convergence theorem, traciality of $\tau$, and $\int_0^\infty (L_p(t) + s)^{-2}\, ds = L_p(t)^{-1}$,
\[
\frac{d}{dt}\tau(\log L_p(t))
= \int_0^\infty \tau\bigl(R_t(s)^2 L'_p(t)\bigr)\, ds
= \tau\bigl(L_p(t)^{-1} L'_p(t)\bigr).
\]
Substituting $L'_p$ and using~\eqref{eq:vn-PQ} and traciality of $\tau$,
\[
\tau(L_p^{-1} L'_p) = 2\tau(H E_p L_p^{-1} E_p^*) + 2\tau(H F_p L_p^{-1} F_p^*) - 2\tau(H) = 2\tau\bigl(H(P_p + Q_p - 1)\bigr).
\]
Summing over $p \in \PP$ yields
\begin{equation}
\label{eq:vn-d1}
\frac{d}{dt}\Phi(e^{tH} g) = 2\tau\bigl(H D(e^{tH} g)\bigr).
\end{equation}

\emph{Step 3: Second derivative and descent.}
Fix $p \in \PP$ and let $V_t := P_p(t) = E_p(t) L_p(t)^{-1}$, which is norm-$C^1$. Using $P_p(t) E_p(t) = E_p(t)$ and $E_p(0) P_p(0) = P_p(0)$, we have two purely algebraic identities:
\begin{align*}
(1 - V_t) e^{tH} V_0 &= (1 - P_p(t)) E_p(t) e^{tH} P_p(0) = 0, \\
e^{tH} V_0 e^{-tH} V_t &= e^{tH} P_p(0) E_p(0) e^{-tH} L_p(t)^{-1} = V_t.
\end{align*}
Differentiating the first identity gives $V'_t e^{tH} V_0 = (1 - V_t) H e^{tH} V_0$; multiplying on the right by $e^{-tH} V_t$ and applying the second identity gives $V'_t V_t = (1 - V_t) H V_t$. Since $V_t^* = V_t$ and $V_t^2 = V_t$, we have $V'_t = V'_t V_t + (V'_t V_t)^* = (1 - V_t) H V_t + V_t H (1 - V_t)$. Therefore, by traciality of $\tau$,
\[
\tau(H V'_t) = \|(1 - V_t) H V_t\|_2^2 + \|V_t H (1 - V_t)\|_2^2 = 2\|(1 - V_t) H V_t\|_2^2.
\]
Because the four components $V_t H V_t$, $V_t H(1 - V_t)$, $(1 - V_t) H V_t$, and $(1 - V_t) H(1 - V_t)$ of $H$ are pairwise orthogonal in $L^2(M, \tau)$, we have $0 \le \tau(H V'_t) \le \|H\|_2^2$. The exact same calculation applies to $Q_p(t)$ (with $F_p$ in place of $E_p$). Summing over all $2n$ projections $P_p, Q_p$ in~\eqref{eq:vn-d1} shows that $t \mapsto \Phi(e^{tH} g)$ is $C^2$ with
\[
0 \le \frac{d^2}{dt^2}\Phi(e^{tH} g) \le 4n \|H\|_2^2.
\]
Taking $H := -D(g)$ and applying Taylor's formula with remainder at $t = \frac{1}{2n}$ yields
\[
\Phi\bigl(e^{-\frac{1}{2n} D(g)} g\bigr)
\le \Phi(g) - 2t \|D(g)\|_2^2 + 2n t^2 \|D(g)\|_2^2
= \Phi(g) - \frac{1}{2n}\|D(g)\|_2^2.
\]
If $\|D(g)\|_2^2 \ge \delta > 0$ held for all $g \in \GL(M)$, iterating $g \mapsto e^{-\frac{1}{2n} D(g)} g$ would decrease $\Phi$ by at least $\frac{\delta}{2n}$ at each step and drive $\Phi$ below $0$, contradicting $\Phi \ge 0$; hence $\inf_{g \in \GL(M)} \|D(g)\|_2 = 0$.

\emph{Step 4: Mixing and meets for bounded idempotents.}
Work in the real inner product space $M_{\mathrm{sa}}$ with inner product $\langle X, Y \rangle_\tau = \tau(XY)$. For any $g \in \GL(M)$, set $A_p := 2P_p(g) - 1$ and $B_{p'} := 2Q_{p'}(g) - 1$; then ${\|A_p\|_2^2 = \|B_{p'}\|_2^2 = 1}$ and $\sum_p A_p + \sum_{l \in \LL} B_{\sigma^{-1}(l)} = 2D(g)$. Applying Lemma~\ref{lem:plane-ineq} (with $\mu = 1$) and expanding $\sum_{(p,p') \in E} \tau(A_p B_{p'}) = 4\sum_{(p,p') \in E} \tau(P_p(g) Q_{p'}(g)) - (q+1)n$ (using $\tau(P_p) + \tau(Q_p) = 1$) gives
\[
4\sum_{(p,p') \in E} \tau\bigl(P_p(g) Q_{p'}(g)\bigr) \le n(q + 1 + \sqrt{q}) + \frac{q + 1}{n}\|D(g)\|_2^2.
\]
For any projections $P, Q \in M$ with meet $W = P \wedge Q$, we have $P(Q - W)P \ge 0$, so $\tau(W) \le \tau(PQ)$. By~(U4), $\tau(r_{f_p} \wedge k_{f_{p'}}) = \tau(P_p(g) \wedge Q_{p'}(g)) \le \tau(P_p(g) Q_{p'}(g))$ for every $g \in \GL(M)$. Taking the infimum over $g \in \GL(M)$ and using $\inf_g \|D(g)\|_2 = 0$ establishes~\eqref{eq:vn-ineq} for all idempotents $f_p \in M$.

\emph{Step 5: Unbounded idempotents.}
Let $e \in \U(M)$ be any idempotent and $P_e := r_e = s_l(e) \in M$. By~(U1), $P_e e = e$ and there exists $c \in \U(M)$ with $ec = P_e$, hence $e P_e = e(ec) = e^2 c = ec = P_e$. Put $a := e - P_e \in \U(M)$; then $P_e a = a$ and $a P_e = 0$, so $a(1 - P_e) = a$ and hence $s_r(a) \le 1 - P_e$ by~(U1). Thus every spectral projection $\1_{(0,k]}(|a|) \le s_r(a) \le 1 - P_e$ commutes with $P_e$, and since $\1_{[0,k]}(|a|) = (1 - s_r(a)) + \1_{(0,k]}(|a|)$ with $P_e \le 1 - s_r(a)$, the projection $\1_{[0,k]}(|a|)$ also commutes with $P_e$ and $1 - P_e$. Setting $a_k := a \1_{[0,k]}(|a|)$, property~(U3) gives $a_k \in M$, and we have $P_e a_k = a_k = a_k(1 - P_e)$, so $a_k \in P_e M (1 - P_e)$ and $a_k^2 = 0$. Consequently, $e_k := P_e + a_k \in M$ is a bounded idempotent satisfying $P_e e_k = e_k$ and $e_k P_e = P_e$, so $e_k \U(M) = P_e\U(M) = e\U(M)$ (its range projection is still $P_e$), while
\[
\rk_\tau(e - e_k) = \rk_\tau\bigl(a\1_{(k,\infty)}(|a|)\bigr) \le \tau\bigl(\1_{(k,\infty)}(|a|)\bigr) \xrightarrow{k \to \infty} 0.
\]

Next, we claim that for any idempotent $h \in \U(M)$ and any projection $w \in M$,
\begin{equation}
\label{eq:vn-inter-id}
w\U(M) \cap (1 - h)\U(M) = \bigl(w - s_r(hw)\bigr)\U(M),
\end{equation}
and hence $\dim_\tau(w\U(M) \cap (1 - h)\U(M)) = \tau(w) - \rk_\tau(hw)$. Indeed, $(hw)w = hw$ implies $s_r(hw) \le w$ by~(U1), so $w_0 := w - s_r(hw)$ is a projection in $M$ with $w_0 \le w$ and $h w_0 = (hw) w_0 = (hw) s_r(hw) w_0 = 0$; thus $w_0 \in w\U(M) \cap (1 - h)\U(M)$. Conversely, if $x \in w\U(M) \cap (1 - h)\U(M)$, then $wx = x$ and $hx = 0$, so $(hw)x = 0$. Choosing $c \in \U(M)$ with $c(hw) = s_r(hw)$ by~(U1) gives $s_r(hw) x = c(hw)x = 0$, hence $w_0 x = wx = x$, proving~\eqref{eq:vn-inter-id}.

Now replace each $f_p \in \U(M)$ by its bounded truncation $f_{p,k} \in M$ constructed above. Since $r_{f_{p,k}} = r_{f_p}$ is independent of $k$, \eqref{eq:vn-inter-id} and~(U2) give
\begin{align*}
\Bigl|\dim_\tau\bigl(f_{p,k}\U(M) &\cap (1 - f_{p',k})\U(M)\bigr)
 - \dim_\tau\bigl(f_p\U(M) \cap (1 - f_{p'})\U(M)\bigr)\Bigr| \\
&\le \rk_\tau\bigl((f_{p',k} - f_{p'}) r_{f_p}\bigr)
\le \rk_\tau(f_{p',k} - f_{p'}) \xrightarrow{k \to \infty} 0.
\end{align*} Passing to the limit $k \to \infty$ in Step~4 applied to $(f_{p,k})_{p \in \PP}$ proves~\eqref{eq:vn-ineq} in full generality.
\end{proof}

\subsection{Proof of Theorem~\ref{thm:vn-main}}

\begin{proof}[Proof of Theorem~\ref{thm:vn-main}]
Let $A$ be any unital ring containing a solution of~\eqref{eq:B-ell} for $\ell = 7$ and some fixed-incidence-free bijection $\sigma$ of $\mathbb{P}^2(\F_7)$ (in particular, $A = S$ by Theorem~\ref{thm:finite-form}(a)).

\emph{Step 1: Unital homomorphisms into $\U(M)$.}
Suppose first that $M$ is a nonzero finite von Neumann algebra admitting a faithful normal tracial state $\tau$, and suppose $\psi\colon A \to \U(M)$ is a unital ring homomorphism. The images of the generators of $B_7^\sigma$ satisfy~\eqref{eq:B-ell} in $\U(M)$. For each $e = (p, p') \in E$, let $w_e \in M$ be the projection with $w_e\U(M) = f_p\U(M) \cap (1 - f_{p'})\U(M)$. Since $f_p x_{e,i} = x_{e,i}$ and $(1 - f_{p'}) x_{e,i} = x_{e,i}$, each $x_{e,i}$ lies in $w_e\U(M)$, so $w_e x_{e,i} = x_{e,i}$ and $\rk_\tau(x_{e,i}) \le \tau(w_e)$ by~(U1). Thus
\[
1 = \rk_\tau(1) = \rk_\tau\Bigl(\sum_{i=0}^2 y_{e,i} x_{e,i}\Bigr) \le \sum_{i=0}^2 \rk_\tau(x_{e,i}) \le 3\tau(w_e),
\]
so $\dim_\tau(f_p\U(M) \cap (1 - f_{p'})\U(M)) \ge \frac{1}{3}$ for all $8n$ edges $e \in E$. By Theorem~\ref{thm:vn-ineq} at $q = 7$,
\[
\frac{8n}{3} \le \sum_{e \in E} \dim_\tau\bigl(f_p\U(M) \cap (1 - f_{p'})\U(M)\bigr) \le \frac{n(8 + \sqrt{7})}{4},
\]
yielding $8 \le 3\sqrt{7}$, i.e., $64 \le 63$, a contradiction.

For an arbitrary nonzero finite von Neumann algebra $M$, let $T\colon M \to Z(M)$ be its center-valued trace~\cite[Theorem~8.2.8]{KR2}. Choose a normal state $\rho$ on the nonzero center $Z(M)$ (for instance, a vector state), and let $z \in Z(M)$ be its support projection ($z \neq 0$). Then $\tau := \rho \circ T|_{zM}$ is a faithful normal tracial state on the nonzero finite von Neumann algebra $zM$: indeed, if $x \in zM_+$ satisfies $\tau(x) = 0$, then $T(x) \in zZ(M)_+$ has $\rho(T(x)) = 0$, hence $T(x) = 0$ by faithfulness of $\rho$ on $zZ(M)$, hence $x = 0$ by faithfulness of $T$. Because $z$ is central, $M = zM \oplus (1 - z)M$ induces a unital ring homomorphism $\U(M) \to \U(zM)$, $x \mapsto x|_{z\HH}$. Composing with any unital homomorphism $A \to \U(M)$ contradicts the previous paragraph.

\emph{Step 2: Non-unital homomorphisms into $\U(M)$ (part~(a)).}
Let $\psi\colon A \to \U(M)$ be a nonzero (not necessarily unital) ring homomorphism. Then $e := \psi(1) \in \U(M)$ is a nonzero idempotent, and $\psi$ is a unital ring homomorphism into the corner $e\U(M)e$. As in Step~5 of the proof of Theorem~\ref{thm:vn-ineq}, let $P_e := r_e \in M$ ($P_e \neq 0$) and $a := e - P_e \in \U(M)$; then $P_e a = a$, $a P_e = 0$, and $a^2 = 0$. Therefore $1 + a$ and $1 - a$ are inverse units in $\U(M)$ satisfying $(1 - a) P_e (1 + a) = (P_e - a P_e)(1 + a) = P_e + P_e a = e$, so $x \mapsto (1 + a)x(1 - a)$ is a unital ring isomorphism $e\U(M)e \xrightarrow{\sim} P_e\U(M)P_e$.

We verify that $P_e\U(M)P_e \cong \U(P_e M P_e)$ as unital rings. If $T \in P_e\U(M)P_e$ has polar decomposition $T = v|T|$, then $P_e T = T = T P_e$ gives $s_l(T) \le P_e$ and $s_r(T) \le P_e$ by~(U1). Thus $v \in P_e M P_e$, every spectral projection $\1_B(|T|)$ for Borel $B \subseteq (0, \infty)$ satisfies $\1_B(|T|) \le s_r(T) \le P_e$ and lies in $P_e M P_e$, and $(1 - P_e)\HH \subseteq \ker T$. Hence $T = T_0 \oplus 0$ on $P_e\HH \oplus (1 - P_e)\HH$, where $T_0 := T|_{P_e\HH}$ is closed and densely defined on $P_e\HH$ with polar partial isometry and spectral projections in $P_e M P_e$, so $T_0 \in \U(P_e M P_e)$. Conversely, every $T_0 \in \U(P_e M P_e)$ yields $T_0 \oplus 0 \in P_e\U(M)P_e$, and $T \mapsto T_0$ preserves strong sums and strong products (closures on common dense domains decompose across $P_e\HH \oplus (1 - P_e)\HH$). Since $P_e M P_e$ is a nonzero finite von Neumann algebra, composing $\psi$ with $e\U(M)e \cong P_e\U(M)P_e \cong \U(P_e M P_e)$ yields a unital ring homomorphism $A \to \U(P_e M P_e)$, contradicting Step~1.

\emph{Step 3: Unital $C^*$-algebras with a tracial state (part~(b)).}
Let $C$ be a unital $C^*$-algebra with a tracial state $\tau$, let $(\pi, \HH_\tau, \xi_\tau)$ be its GNS representation, and set $N := \pi(C)'' \subseteq B(\HH_\tau)$ and $\varphi(x) := \langle x\xi_\tau, \xi_\tau \rangle$ for $x \in N$. Then $\varphi$ is a normal state on $N$ (in particular $N \neq 0$) that is tracial on $\pi(C)$. Since $\pi(C)$ is weak-operator dense in $N$ by the bicommutant theorem, for any fixed $b \in \pi(C)$ and $x \in N$, taking a net $a_\alpha \in \pi(C)$ converging to $x$ in the weak operator topology and using traciality of $\varphi$ on $\pi(C)$ gives
\[
\varphi(xb)
= \langle x(b\xi_\tau), \xi_\tau \rangle
= \lim_\alpha \langle a_\alpha b\xi_\tau, \xi_\tau \rangle
= \lim_\alpha \varphi(a_\alpha b)
= \lim_\alpha \varphi(b a_\alpha)
= \lim_\alpha \langle a_\alpha \xi_\tau, b^*\xi_\tau \rangle
= \varphi(bx).
\]
Approximating an arbitrary $y \in N$ in the weak operator topology by a net $b_\beta \in \pi(C)$ then gives $\varphi(xy) = \lim_\beta \varphi(x b_\beta) = \lim_\beta \varphi(b_\beta x) = \varphi(yx)$, so $\varphi$ is tracial on $N$. If $x \in N$ satisfies $\varphi(x^*x) = 0$, then $x\xi_\tau = 0$, and for all $a, b \in \pi(C)$ we have $\langle x a\xi_\tau, b\xi_\tau \rangle = \varphi(b^* x a) = \varphi(a b^* x) = \langle x\xi_\tau, b a^* \xi_\tau \rangle = 0$; since $\pi(C)\xi_\tau$ is dense in $\HH_\tau$, $x = 0$. Thus $N$ is a nonzero finite von Neumann algebra, and any unital ring homomorphism $A \to C$ composes with $\pi\colon C \to N \subseteq \U(N)$ to contradict Step~1.

\emph{Step 4: Non-unital homomorphisms into rank ultraproducts (part~(c)).}
Let $\psi\colon A \to \mathcal{Q} := \RUF{\omega}{d_i}{F_i}$ be a nonzero ring homomorphism with $\operatorname{char}(F_i) = 0$ and $d_i \ge 1$. Put $e := \psi(1) \neq 0$, and lift $e$ to an exact idempotent sequence $(E_i)$ in $\prod_{i \in I} M_{d_i}(F_i)$ via Lemma~\ref{lem:lifting}(a). Since $e \neq 0$, there exist $c > 0$ and a set $U \in \omega$ such that $r_i := \rank(E_i) \ge c d_i > 0$ for all $i \in U$. For each $i \in U$, choosing bases adapted to $\im(E_i) \oplus \ker(E_i)$ identifies the corner $E_i M_{d_i}(F_i) E_i$ with $M_{r_i}(F_i)$ while preserving matrix rank. Because $c \le r_i / d_i \le 1$ on $U$, letting $\omega|_U := \{V \subseteq U : V \in \omega\}$ be the restricted ultrafilter on $U$, a sequence $(X_i)_{i \in U}$ with $X_i \in E_i M_{d_i}(F_i) E_i$ satisfies $\lim_{i \to \omega} \rank(X_i)/d_i = 0$ if and only if $\lim_{i \to \omega|_U} \rank(X_i)/r_i = 0$. Consequently, the unital corner $e\mathcal{Q}e$ is isomorphic to the rank ultraproduct $\mathcal{Q}' := \RUF{\omega|_U}{r_i}{F_i}$ with $r_i \ge 1$. Thus $\psi$ yields a unital ring homomorphism $A \to \mathcal{Q}'$, carrying the solution of~\eqref{eq:B-ell} in $A$ into the von Neumann regular ring $\mathcal{Q}'$. By Lemma~\ref{lem:B-equiv}(b), $\mathcal{Q}'$ then contains a solution of~\eqref{eq:B-tilde}, contradicting Proposition~\ref{prop:no-config}.
\end{proof}

\section{Formal verification in Lean~4 and open problems}
\label{sec:formal}

\subsection{Formalization map and verification record}

The theorems and proofs of Sections~\ref{sec:ineq}--\ref{sec:proofs} have been formalized in Lean~4 (toolchain \code{v4.34.1}, Mathlib commit \code{d13f23b}) and are publicly available at \url{https://github.com/jackryan225/rank-rings-char0}. A standalone specification file \code{Challenge/Statement.lean} (importing Mathlib only) defines \code{IsVNRegular}, \code{IsRankFunction}, \code{IsRankNull}, and \code{RankUltraproduct}, and states the three headline theorems \code{elek\_countable}, \code{elek\_complex}, and \code{finite\_form}. The file \code{Solution.lean} proves those exact declarations by transferring from the modular development under \code{RankRings/}. As noted in Section~\ref{sec:intro}, \code{Challenge/Statement.lean} states \code{finite\_form} in existential form over $\C$; the explicit $B_7^\sigma$, all-characteristic-zero-fields, and characteristic-$7$ contrast versions of Theorem~\ref{thm:finite-form} are proved as internal theorems in \code{FiniteFormF.lean} and \code{ModL.lean} (built and axiom-checked, though not compared by \code{Check/SameStatement.lean}), quantifying over fields in universe $0$ (\code{Type}), which is mathematically harmless since any homomorphism from a countable ring or matrix assignment for a finite presentation factors through countable subfields and matrix rank is invariant under field extension. Table~\ref{tab:lean-map} summarizes the correspondence between the mathematical results in this paper and their Lean~4 formalizations.

\begin{table}[ht]
\centering
\caption{Correspondence between paper results and Lean~4 declarations (modules lie in \code{RankRings/} unless prefixed otherwise; all filenames carry the \code{.lean} suffix).}
\label{tab:lean-map}
\footnotesize
\setlength{\tabcolsep}{5pt}
\begin{tabular}{@{}ll>{\raggedright\arraybackslash}p{2.9in}@{}}
\toprule
Paper result & Lean module(s) & Primary declaration(s) \\
\midrule
Lemmas~\ref{lem:ld-trace}, \ref{lem:step-ineq} & \texttt{LogDet} & \texttt{ld\_le\_trace}, \texttt{ld\_add\_ld\_le}, \texttt{step\_ineq} \\
Lemma~\ref{lem:L-idem}, Prop.~\ref{prop:balanced} & \texttt{Descent} & \texttt{comparison}, \texttt{comparison\_step}, \texttt{exists\_balanced} \\
Lemma~\ref{lem:plane-ineq} & \texttt{Plane} & \texttt{plane\_ineq} \\
Theorem~\ref{thm:ineq-C} & \texttt{Ineq} & \texttt{ineq\_C}, \texttt{trace\_mul\_ge} \\
Lemma~\ref{lem:srk-basic}, Prop.~\ref{prop:smith} & \texttt{SRank}, \texttt{Smith} & \texttt{srk\_mul\_le}, \texttt{hasSmith\_all}, \texttt{exists\_qinv} \\
Theorem~\ref{thm:RR-regular} & \texttt{RankRing}, \texttt{RankFn}, \texttt{Main} & \texttt{RR}, \texttt{rk\_add\_orth}, \texttt{rr\_regular}, \texttt{isRankFunction\_RR} \\
Lemma~\ref{lem:hall-sigma} & \texttt{Geom} & \texttt{exists\_sigma}, \texttt{exists\_lift} \\
Lemmas~\ref{lem:exact-third}, \ref{lem:offblock-anti} & \texttt{Words}, \texttt{WordsBlock} & \texttt{sum\_transpose\_M\_mul\_M}, \texttt{offblock}, \texttt{window\_le} \\
Theorem~\ref{thm:exotic} & \texttt{Exotic}, \texttt{ExoticS} & \texttt{exotic}, \texttt{isNull\_of\_approx}, \texttt{exotic\_RR} \\
Proposition~\ref{prop:modL} & \texttt{ModL} & \texttt{offblock\_zero}, \texttt{exists\_config\_modL}, \texttt{rel\_modL} \\
Lemma~\ref{lem:lifting}, Prop.~\ref{prop:no-config} & \texttt{NoHom}, \texttt{NoHomF} & \texttt{Gen.idem\_lift}, \texttt{Gen.exists\_embed}, \texttt{Gen.no\_config} \\
$\overline{\Q}_7 \cong \C$ (choice) & \texttt{CAlg} & \texttt{nonempty\_padicAlgCl\_equiv\_complex} \\
Theorem~\ref{thm:main-ultra} (checked) & \texttt{Main}, \texttt{Final}, \texttt{Solution} & \texttt{main\_core}, \texttt{elek\_countable}, \texttt{elek\_complex} \\
Thm.~\ref{thm:finite-form} (checked, $\C$) & \texttt{FiniteForm}, \texttt{Solution} & \texttt{finite\_form} \\
Lemma~\ref{lem:B-equiv}, Thm.~\ref{thm:finite-form} (internal) & \texttt{FiniteFormF}, \texttt{ModL} & \texttt{relP\_of\_config}, \texttt{finite\_form\_paper}, \texttt{eps\_char0\_paper}, \texttt{char\_contrast} \\
Definition sanity checks & \texttt{Check/Sanity} & \texttt{quotient\_eq}, \texttt{ultra\_ne\_zero}, \texttt{rat\_maps} \\
\bottomrule
\end{tabular}
\end{table}

On November 25, 2026, the repository was built and audited from a clean checkout with the following outcomes (logs recorded in \code{logs/}):
\begin{itemize}[leftmargin=1.8em]
\item \code{lake build} succeeded across 8,948 jobs with zero occurrences of \code{sorry}, \code{admit}, \code{axiom}, or \code{native\_decide} in \code{RankRings/}, \code{RankRings.lean}, and \code{Solution.lean}.
\item \code{Check/Axioms.lean} confirmed that all 21 exported theorems depend exclusively on the three standard Lean axioms \code{propext}, \code{Classical.choice}, and \code{Quot.sound}.
\item \code{Check/SameStatement.lean} verified that \code{Challenge/Statement.lean} imports Mathlib only and that all 20 of its declarations are definitionally and syntactically identical to those in \code{Solution.lean} (with byte-identical \code{pp.all} views).
\item \code{Check/Sanity.lean} verified non-degeneracy of the definitions: \code{RankUltraproduct} agrees with the quotient by $J_\omega$ (\code{quotient\_eq}), is nonzero whenever all $d_i > 0$ (\code{ultra\_ne\_zero}), and receives unital homomorphisms from $\Q$, $\C$, and $M_2(\Q)$ (\code{rat\_maps}, \code{complex\_maps}, \code{matrix\_maps}).
\item Independent kernel verification succeeded with both \code{leanchecker} (replaying all 23 modules: \code{RankRings}, its 21 submodules, and \code{Solution}) and \code{lean4lean} (checking all 68,066 declarations, including 53,498 theorems, in the export of \code{elek\_countable}, \code{elek\_complex}, and \code{finite\_form}).
\end{itemize}
As noted above, Proposition~\ref{prop:stably-finite} and Section~\ref{sec:vn} (Theorems~\ref{thm:vn-main} and~\ref{thm:vn-ineq}) are paper proofs and have not been formalized in Lean~4.

\subsection{Open problems}

We close with four natural questions left open by the present construction.

\begin{question}
Does there exist a regular rank $\C$-algebra of \emph{countable dimension over $\C$} that admits no unital $\C$-algebra embedding (or no unital ring homomorphism) into Elek's matrix rank ultraproducts $M^{\mathrm{alg}}_\nu$?
\end{question}

\begin{question}
Does our countable regular rank $\Q$-algebra $S$ admit a unital ring homomorphism into a rank ultraproduct $\RUF{\omega}{d_i}{\F_{p_i}}$ of matrix rings over finite fields of growing prime characteristics $p_i \to \infty$? More generally, is there a countable regular rank $\Q$-algebra admitting no unital ring homomorphism into \emph{any} rank ultraproduct of matrix rings over fields of arbitrary characteristics?
\end{question}

\begin{question}
Does the incidence inequality~\eqref{eq:vn-ineq} hold in the regular ring of a finite $AW^*$-algebra that is not a von Neumann algebra (where a normal scalar-valued trace is not available)?
\end{question}

\begin{question}
Can the countable regular rank $\Q$-algebra $S$ of Theorem~\ref{thm:main-ultra}\textup{(a)} be chosen to be simple, or to have a unique rank function?
\end{question}

\section*{Acknowledgments and disclosure of AI assistance}

This work was carried out by an independent high school author in collaboration with large language models (Claude Opus~5.5, Claude Fable~5.2, GPT-6.2 Astra, and Gemini~4 Neon), which assisted in discovering the proof architecture, drafting and auditing the Lean~4 formalization, checking the paper proofs in Section~\ref{sec:vn}, and preparing the manuscript. All results in Sections~\ref{sec:ineq}--\ref{sec:proofs} are mechanically verified by the Lean~4 kernel and independent checkers as documented in Section~\ref{sec:formal}; the non-unital and operator-algebraic arguments in Section~\ref{sec:vn} (Theorems~\ref{thm:vn-main} and~\ref{thm:vn-ineq}) and Proposition~\ref{prop:stably-finite} are paper proofs and have not yet undergone formal refereeing. The author warmly thanks G\'abor Elek for encouraging this write-up and for his clear framing of the problem in~\cite{Elek2013}.

\begin{thebibliography}{99}

\bibitem{AP2017}
G.~Arzhantseva and L.~P\u{a}unescu,
\emph{Linear sofic groups and algebras},
Trans.\ Amer.\ Math.\ Soc.\ \textbf{369} (2017), no.~4, 2285--2310,
\href{https://arxiv.org/abs/1212.6780}{arXiv:1212.6780}.

\bibitem{BDWY2011}
B.~Barak, Z.~Dvir, A.~Wigderson, and A.~Yehudayoff,
\emph{Rank bounds for design matrices with applications to combinatorial geometry and locally correctable codes},
Proc.\ 43rd ACM Symposium on Theory of Computing (STOC~2011), 519--528,
\href{https://arxiv.org/abs/1009.4375}{arXiv:1009.4375}.

\bibitem{Barthe1998}
F.~Barthe,
\emph{On a reverse form of the Brascamp--Lieb inequality},
Invent.\ Math.\ \textbf{134} (1998), no.~2, 335--361.

\bibitem{Berberian1957}
S.~K.~Berberian,
\emph{The regular ring of a finite $AW^*$-algebra},
Ann.\ of Math.\ (2) \textbf{65} (1957), no.~2, 224--240.

\bibitem{DvirHu2016}
Z.~Dvir and G.~Hu,
\emph{Sylvester--Gallai for arrangements of subspaces},
Discrete Comput.\ Geom.\ \textbf{56} (2016), no.~4, 940--965,
\href{https://arxiv.org/abs/1412.0795}{arXiv:1412.0795}.

\bibitem{Elek2013}
G.~Elek,
\emph{Connes embeddings and von Neumann regular closures of group algebras},
Trans.\ Amer.\ Math.\ Soc.\ \textbf{365} (2013), no.~6, 3019--3039,
\href{https://arxiv.org/abs/1006.5378}{arXiv:1006.5378}.

\bibitem{ElekSzabo}
G.~Elek and E.~Szab\'o,
\emph{Sofic groups and direct finiteness},
J.\ Algebra \textbf{280} (2004), no.~2, 426--434,
\href{https://arxiv.org/abs/math/0305440}{arXiv:math/0305440}.

\bibitem{Goodearl}
K.~R.~Goodearl,
\emph{Von Neumann Regular Rings},
2nd ed., Robert E.\ Krieger Publishing Co., Malabar, FL, 1991.

\bibitem{Greenfeld2023}
B.~Greenfeld,
\emph{Stable finiteness does not imply linear soficity},
Math.\ Proc.\ Cambridge Philos.\ Soc.\ \textbf{175} (2023), no.~2, 319--325,
\href{https://arxiv.org/abs/2210.11650}{arXiv:2210.11650}.

\bibitem{KR2}
R.~V.~Kadison and J.~R.~Ringrose,
\emph{Fundamentals of the Theory of Operator Algebras, Volume~II: Advanced Theory},
Academic Press, Orlando, FL, 1986.

\bibitem{Luck2002}
W.~L\"uck,
\emph{$L^2$-Invariants: Theory and Applications to Geometry and $K$-Theory},
Ergebnisse der Mathematik und ihrer Grenzgebiete (3), vol.~44, Springer-Verlag, Berlin, 2002.

\bibitem{MurrayVonNeumann}
F.~J.~Murray and J.~von Neumann,
\emph{On rings of operators},
Ann.\ of Math.\ (2) \textbf{37} (1936), no.~1, 116--229.

\bibitem{Schofield1985}
A.~H.~Schofield,
\emph{Representations of Rings over Skew Fields},
London Mathematical Society Lecture Note Series, vol.~92, Cambridge University Press, Cambridge, 1985.

\end{thebibliography}

\end{document}
